\documentclass[12pt]{article}
\usepackage[T1]{fontenc}
\usepackage[utf8]{inputenc}
\usepackage{amsmath,amssymb,booktabs,longtable,array,geometry,microtype,enumitem}
\geometry{letterpaper,margin=0.92in}
\setlength{\parindent}{1.4em}
\setlength{\parskip}{0.35em}
\setlength{\emergencystretch}{2em}
\allowdisplaybreaks
\newcommand{\art}[1]{\subsection*{Art. #1}}
\newcommand{\norm}[1]{\operatorname{N}(#1)}
\newcommand{\gbrack}[1]{\left[#1\right]}
\newcolumntype{L}{>{$}l<{$}}
\newcolumntype{C}{>{$}c<{$}}
\sloppy

\begin{document}
\begin{center}
{\Large C. F. Gauss}\par
{\large Theoria residuorum biquadraticorum. Commentatio secunda}\par
{\large Articuli 30--76}
\end{center}

\art{30}

Amplissimam itaque messem theorematum specialium aperit inductio, theoremati pro numero $2$ affinium: sed desideratur vinculum commune, desiderantur demonstrationes rigorosae, quum methodus, per quam in commentatione prima numerum $2$ absolvimus, ulteriorem applicationem non patiatur. Non desunt quidem methodi diversae, per quas demonstrationibus pro casibus particularibus potiri liceret, iis potissimum, qui distributionem residuorum quadraticorum inter complexus $A$, $C$ spectant, quibus tamen non immoramur; quum theoria generalis omnes casus complectens in votis esse debeat.

Cui rei quum inde ab anno 1805 meditationes nostras dicare coepissemus, mox certiores facti sumus, fontem genuinum theoriae generalis in campo arithmeticae promoto quaerendum esse, uti iam in art.~1 addigitavimus.

Quemadmodum scilicet arithmetica sublimior in quaestionibus hactenus pertractatis inter solos numeros integros reales versatur, ita theoremata circa residua biquadratica tunc tantum in summa simplicitate ac genuina venustate resplendent, quando campus arithmeticae ad quantitates imaginarias extenditur, ita ut absque restrictione ipsius obiectum constituant numeri formae $a+bi$, denotantibus $i$, pro more quantitatem imaginariam $\sqrt{-1}$, atque $a$, $b$ indefinite omnes numeros reales integros inter $-\infty$ et $+\infty$. Tales numeros vocabimus \textit{numeros integros complexos}, ita quidem, ut reales complexis non opponantur, sed tamquam species sub his contineri censeantur. Commentatio praesens tum doctrinam elementarem de numeris complexis, tum prima initia theoriae residuorum biquadraticorum sistet, quam ab omni parte perfectam reddere in continuatione subsequente suscipiemus\footnote{Obiter saltem hic adhuc monere convenit, campum ita definitum imprimis theoriae residuorum biquadraticorum accommodatum esse. Theoria residuorum cubicorum simili modo superstruenda est considerationi numerorum formae $a+bh$, ubi $h$ est radix imaginaria aequationis $h^3-1=0$, puta $h=-\frac{1}{2}+\sqrt{\frac{3}{4}}\,i$; et perinde theoria residuorum potestatum altiorum introductionem aliarum quantitatum imaginariarum postulabit.}.

\art{31}

Ante omnia quasdam denominationes praemittimus, per quarum introductionem brevitati et perspicuitati consuletur. Campus numerorum complexorum $a+bi$ continet

\begin{enumerate}[label=\Roman*.]
\item numeros reales, ubi $b=0$, et, inter hos, pro indole ipsius $a$,
    \begin{enumerate}[label=\arabic*)]
    \item cifram,
    \item numeros positivos,
    \item numeros negativos;
    \end{enumerate}
\item numeros imaginarios, ubi $b$ cifrae inaequalis. Hic iterum distinguuntur
    \begin{enumerate}[label=\arabic*)]
    \item numeri imaginarii absque parte reali, i.e. ubi $a=0$,
    \item numeri imaginarii cum parte reali, ubi neque $b$ neque $a=0$.
    \end{enumerate}
\end{enumerate}

Priores si placet \textit{numeri imaginarii puri}, posteriores \textit{numeri imaginarii mixti} vocari possunt.

Unitatibus in hac doctrina utimur quaternis, $+1$, $-1$, $+i$, $-i$, quae simpliciter positiva, negativa, positiva imaginaria, negativa imaginaria audient. Producta terna cuiuslibet numeri complexi per $-1$, $+i$, $-i$ illius \textit{socios} vel numeros \textit{illi associatos} appellabimus. Excepta itaque cifra (quae sibi ipsa associata est), semper quaterni numeri inaequales associati sunt.

Contra numero complexo \textit{coniunctum} vocamus eum, qui per permutationem ipsius $i$ cum $-i$ inde oritur. Inter numeros imaginarios itaque bini inaequales semper coniuncti sunt, dum numeri reales sibi ipsi sunt coniuncti, siquidem denominationem ad hos extendere placet.

Productum numeri complexi per numerum ipsi coniunctum utriusque \textit{normam} vocamus. Pro norma itaque numeri realis, ipsius quadratum habendum est. Generaliter octonos numeros nexos habemus, puta
\[
\begin{array}{r|r}
 a+bi & a-bi\\
 -b+ai & -b-ai\\
 -a-bi & -a+bi\\
 b-ai & b+ai
\end{array}
\]
ubi duas quaterniones numerorum associatorum, quatuor biniones coniunctorum conspicimus, omniumque norma communis est $aa+bb$. Sed octo numeri ad quatuor inaequales reducuntur, quoties vel $a=\pm b$, vel alteruter numerorum $a$, $b=0$.

E definitionibus allatis protinus demanant sequentia. Producto duorum numerorum complexorum coniunctum est productum e numeris, qui illis coniuncti sunt. Idem valet de producto e pluribus factoribus, nec non de quotientibus. Norma producti e duobus numeris complexis aequalis est producto ex horum normis. Hoc quoque theorema extenditur ad producta e quotcunque factoribus et ad quotientes. Cuiusvis numeri complexi (excipiendo cifram, quod plerumque abhinc tacite subintelligemus) norma est numerus positivus.

Ceterum nihil obstat, quominus definitiones nostrae ad valores fractos vel adeo irrationales ipsorum $a$, $b$ extendantur; sed $a+bi$ tunc tantum \textit{numerus complexus integer} audiet, quando uterque $a$, $b$ est integer, atque tunc tantum \textit{rationalis}, quando uterque $a$, $b$ rationalis est.

\art{32}

Algorithmus operationum arithmeticarum circa numeros complexos vulgo notus est: divisio, per introductionem normae, ad multiplicationem reducitur, quum habeatur
\[
\frac{a+bi}{c+di}=(a+bi)\cdot\frac{c-di}{cc+dd}=\frac{ac+bd}{cc+dd}+\frac{bc-ad}{cc+dd}\,i.
\]
Extractio radicis quadratae perficitur adiumento formulae
\[
\sqrt{a+bi}=\pm\left(\sqrt{\frac{\sqrt{aa+bb}+a}{2}}+i\sqrt{\frac{\sqrt{aa+bb}-a}{2}}\right)
\]
si $b$ est numerus positivus, vel huius
\[
\sqrt{a+bi}=\pm\left(\sqrt{\frac{\sqrt{aa+bb}+a}{2}}-i\sqrt{\frac{\sqrt{aa+bb}-a}{2}}\right)
\]
si $b$ est numerus negativus. Usui transformationis quantitatis complexae $a+bi$ in $r(\cos \varphi+i\sin\varphi)$ ad calculos facilitandos, non opus est hic immorari.

\art{33}

Numerum integrum complexum, qui in factores duos ab unitatibus diversos resolvi potest\footnote{Sive, quod idem est, tales, quorum normae unitate sint maiores.}, vocamus \textit{numerum complexum compositum}; contra \textit{numerus primus complexus} dicetur, qui talem resolutionem in factores non admittit. Hinc statim patet, quemvis numerum compositum realem etiam esse compositum complexum.

At numerus primus realis poterit esse numerus complexus compositus, et quidem hoc valebit de numero $2$ atque de omnibus numeris primis realibus positivis formae $4n+1$ (excepto numero 1), quippe quos in bina quadrata positiva decomponi posse constat; puta, fit
\[
2=(1+i)(1-i),\quad 5=(1+2i)(1-2i),\quad 13=(3+2i)(3-2i),\quad 17=(1+4i)(1-4i),\quad \text{etc.}
\]
Contra numeri primi reales positivi formae $4n+3$ semper sunt numeri primi complexi. Si enim talis numerus $q$ esset $(a+bi)(\alpha+\beta i)$, foret etiam $q=(a-bi)(\alpha-\beta i)$, adeoque
\[
qq=(aa+bb)(\alpha\alpha+\beta\beta):
\]
at $qq$ unico tantum modo in factores positivos unitate maiores resolvi potest, puta in $q\times q$, unde esse deberet $q=aa+bb=\alpha\alpha+\beta\beta$, Q.E.A.; quum summa duorum quadratorum nequeat esse formae $4n+3$.

Numeri reales negativi manifesto easdem denominationes servant, quas positivi, idemque valet de numeris imaginariis puris.

Superest itaque, ut inter numeros imaginarios mixtos, compositos a primis dignoscere doceamus, quod fit per sequens

\medskip
\noindent\textsc{Theorema.} Quivis numerus integer imaginarius mixtus $a+bi$ est vel numerus primus complexus, vel numerus compositus, prout ipsius norma est vel numerus primus realis, vel numerus compositus.

\medskip
\noindent\textsc{Dem.} I. Quoniam numeri complexi compositi norma semper est numerus compositus, patet, numerum complexum, cuius norma sit numerus primus realis, necessario esse debere numerum primum complexum. Q.E.P.

II. Si vero norma $aa+bb$ est numerus compositus, sit $p$ numerus primus positivus realis illam metiens. Duo iam casus distinguendi sunt.

1) Si $p$ est formae $4n+3$, constat, $aa+bb$ per $p$ divisibilem esse non posse, nisi $p$ simul metiatur ipsos $a$, $b$, unde $a+bi$ erit numerus compositus.

2) Si $p$ non est formae $4n+3$, certo in duo quadrata decomponi poterit: statuemus itaque $p=\alpha\alpha+\beta\beta$. Quum fiat
\[
(a\alpha+b\beta)(a\alpha-b\beta)=aa(\alpha\alpha+\beta\beta)-\beta\beta(aa+bb)
\]
adeoque per $p$ divisibilis, $p$ certo alterutrum factorem $a\alpha+b\beta$, $a\alpha-b\beta$ metietur, et quum insuper fiat
\[
(a\alpha+b\beta)^2+(b\alpha-a\beta)^2=(a\alpha-b\beta)^2+(b\alpha+a\beta)^2=(aa+bb)(\alpha\alpha+\beta\beta)
\]
adeoque per $pp$ divisibilis, patet, in casu priori etiam $b\alpha-a\beta$, in posteriori $b\alpha+a\beta$ per $p$ divisibilem esse debere. Quare in casu priori
\[
\frac{a+bi}{\alpha+\beta i}=\frac{a\alpha+b\beta}{p}+\frac{b\alpha-a\beta}{p}\,i
\]
erit numerus integer complexus, in posteriori autem
\[
\frac{a+bi}{\alpha-\beta i}=\frac{a\alpha-b\beta}{p}+\frac{b\alpha+a\beta}{p}\,i
\]
integer erit. Quum itaque numerus propositus vel per $\alpha+\beta i$ vel per $\alpha-\beta i$ divisibilis sit, quotientisque norma $=\frac{aa+bb}{p}$ per hyp. ab unitate diversa fiat, patet, $a+bi$ in utroque casu esse numerum complexum compositum. Q.E.S.

\art{34}

Totum itaque ambitum numerorum primorum complexorum exhauriunt quatuor species sequentes:
\begin{enumerate}[label=\arabic*)]
\item quatuor unitates, $1$, $+i$, $-1$, $-i$, quas tamen, dum de numeris primis agemus, plerumque tacite subintelligemus exclusas;
\item numerus $1+i$ cum tribus sociis $-1+i$, $-1-i$, $1-i$;
\item numeri primi reales positivi formae $4n+3$ cum ternis sociis;
\item numeri complexi, quorum normae sunt numeri primi reales formae $4n+1$ unitate maiores, et quidem cuivis normae tali datae semper octoni numeri primi complexi et non plures respondebunt, quum talis norma unico tantum modo in bina quadrata decomponi possit.
\end{enumerate}

\art{35}

Quemadmodum numeri integri reales in pares et impares distribuuntur, atque illi iterum in pariter pares et impariter pares, ita inter numeros complexos distinctio aeque essentialis se offert: sunt scilicet

\textit{vel} per $1+i$ non divisibiles, puta numeri $a+bi$, ubi alter numerorum $a$, $b$ est impar, alter par;

\textit{vel} per $1+i$ neque vero per $2$ divisibiles, quoties uterque $a$, $b$ est impar;

\textit{vel} per $2$ divisibiles, quoties uterque $a$, $b$ est par.

Numeri primae classis commode dici possunt \textit{numeri complexi impares}, secundae \textit{semipares}, tertiae \textit{pares}.

Productum e pluribus factoribus complexis semper impar erit, quoties omnes factores sunt impares; semipar, quoties unus factor est semipar, reliqui impares; par autem, quoties inter factores vel saltem duo semipares inveniuntur, vel saltem unus par.

Norma cuiusvis numeri complexi imparis est formae $4n+1$; norma numeri semiparis est formae $8n+2$; denique norma numeri paris est productum numeri formae $4n+1$ in numerum $4$ vel altiorem binarii potestatem.

\art{36}

Quum nexus inter quaternos numeros complexos socios analogus sit nexui inter binos numeros reales oppositos (i.e. absolute aequales signisque oppositis affectos), atque ex his vulgo positivus tamquam primarius merito considerari soleat: quaestio oritur, num similis distinctio inter quaternos numeros complexos socios stabiliri possit, et pro utili haberi debeat.

Ad quam decidendam perpendere oportet, principium distinctionis ita comparatum esse debere, ut productum duorum numerorum, qui inter socios suos pro primariis valent, semper fiat numerus primarius inter socios suos. At mox certiores fimus, tale principium omnino non dari, nisi distinctio ad numeros integros restringatur: quinadeo distinctio utilis ad numeros impares limitanda erit. Pro his vero finis propositus duplici modo attingi potest. Scilicet

I. Productum duorum numerorum $a+bi$, $a'+b'i$ ita comparatorum, ut $a$, $a'$ sint formae $4n+1$, atque $b$, $b'$ pares, eadem proprietate gaudebit, ut pars realis fiat $\equiv 1 \pmod{4}$, atque pars imaginaria par. Et facile perspicietur, inter quaternos numeros impares associatos unum solum sub illa forma contentum esse.

II. Si numerus $a+bi$ ita comparatus est, ut $a-1$ et $b$ vel simul pariter pares sint, vel simul impariter pares, eius productum per numerum complexum eiusdem formae eadem forma gaudebit, facileque perspicitur, e quaternis numeris imparibus associatis unum solum sub hac forma contineri.

Ex his duobus principiis aeque fere idoneis posterius adoptabimus, scilicet inter quaternos numeros complexos impares associatos eum pro primario habebimus, qui secundum modulum $2+2i$ unitati positivae fit congruus: hoc pacto plura insignia theoremata maiori concinnitate enunciare licebit. Ita e.g. sunt numeri primi complexi primarii $-1+2i$, $-1-2i$, $+3+2i$, $+3-2i$, $+1+4i$, $+1-4i$ etc., nec non reales $-3$, $-7$, $-11$, $-19$ etc. manifesto semper signo negativo afficiendi. Numero complexo impari primario coniunctus quoque primarius erit.

Pro numeris semiparibus et paribus in genere similis distinctio nimis arbitraria parumque utilis foret. E numeris primis associatis $1+i$, $1-i$, $-1+i$, $-1-i$ unum quidem prae reliquis pro primario eligere possumus, sed ad compositos talem distinctionem non extendemus.

\art{37}

Si inter factores numeri complexi compositi inveniuntur tales, qui ipsi sunt compositi, atque hi iterum in factores suos resolvuntur, manifesto tandem ad factores primos delabimur, i.e. quivis numerus compositus in factores primos resolubilis est. Inter quos si qui non primarii reperiuntur, singulorum loco substituatur productum primarii associati per $i$, $-1$ vel $-i$. Hoc pacto patet, quemvis numerum complexum compositum $M$ reduci posse ad formam
\[
M=i^{\mu}A^{\alpha}B^{\beta}C^{\gamma}\ldots
\]
ita ut $A$, $B$, $C$ etc. sint numeri primi complexi primarii inaequales, atque $\mu=0$, $1$, $2$ vel $3$. Circa hanc resolutionem theorema se offert, unico tantum modo eam fieri posse, quod theorema obiter quidem consideratum per se manifestum videri posset, sed utique demonstratione eget. Ad quam sternit viam sequens

\medskip
\noindent\textsc{Theorema.} Productum $M=A^{\alpha}B^{\beta}C^{\gamma}\ldots$, denotantibus $A$, $B$, $C$ etc. numeros primos complexos primarios diversos, divisibile esse nequit per ullum numerum primum complexum primarium, qui inter $A$, $B$, $C$ etc. non reperitur.

\medskip
\noindent\textsc{Dem.} Sit $P$ numerus primus complexus primarius inter $A$, $B$, $C$ etc. non contentus, sintque $p$, $a$, $b$, $c$ etc. normae numerorum $P$, $A$, $B$, $C$ etc. Hinc facile colligitur, normam numeri $M$ fore $=a^{\alpha}b^{\beta}c^{\gamma}$ etc., unde hic numerus, si $M$ per $P$ divisibilis esset, per $p$ divisibilis esse deberet.

Quum singulae normae sint vel numeri primi reales (e serie $2$, $5$, $13$, $17$ etc.), vel numerorum primorum realium quadrata (e serie $9$, $49$, $121$ etc.), sponte patet, illud evenire non posse, nisi $p$ cum aliqua norma $a$, $b$, $c$ etc. identica fiat: supponemus itaque $p=a$. At quum $P$, $A$ per hyp. sint numeri primi complexi primarii non identici, facile perspicietur, haec simul consistere non posse, nisi $P$, $A$ sint numeri complexi imaginarii coniuncti, et proin $p=a$ numerus primus realis impar (non quadratum numeri primi): supponemus itaque $A=k+li$, $P=k-li$.

Hinc (extendendo notionem et signum congruentiae ad numeros integros complexos) erit $A\equiv 2k \pmod{P}$, unde facile colligitur
\[
M\equiv 2^{\alpha}k^{\alpha}B^{\beta}C^{\gamma}\ldots \pmod{P}.
\]
Quapropter dum $M$ per $P$ divisibilis supponitur, erit etiam $2^{\alpha}k^{\alpha}B^{\beta}C^{\gamma}\ldots$ per $P$ divisibilis, adeoque norma huius numeri, quae fit $=2^{2\alpha}k^{2\alpha}b^{\beta}c^{\gamma}\ldots$, divisibilis per $p$. At quum $2$ et $k$ per $p$ certo non sint divisibiles, hinc sequitur, $p$ cum aliquo numerorum $b$, $c$ etc. identicum esse debere: sit e.g. $p=b$. Hinc vero concludimus, esse vel $B=k+li$, vel $B=k-li$, i.e. vel $B=A$, vel $B=P$, utrumque contra hyp.

Ex hoc theoremate alterum, quod resolutio in factores primos unico tantum modo perfici potest, facillime derivatur, et quidem per ratiocinia iis, quibus in \textit{Disquisitionibus Arithmeticis} pro numeris realibus usi sumus (art.~16), prorsus analoga: quapropter illis hic immorari superfluum foret.

\art{38}

Progredimur iam ad congruentiam numerorum secundum modulos complexos. Sed in limine huius disquisitionis convenit indicare, quomodo ditio quantitatum complexarum intuitui subiici possit.

Sicuti omnis quantitas realis per partem rectae utrinque infinitae ab initio arbitrario sumendam, et secundum segmentum arbitrarium pro unitate acceptum aestimandam exprimi, adeoque per punctum alterum repraesentari potest, ita ut puncta ab altera initii plaga quantitates positivas, ab altera negativas repraesentent: ita quaevis quantitas complexa repraesentari poterit per aliquod punctum in plano infinito, in quo recta determinata ad quantitates reales refertur, scilicet quantitas complexa $x+iy$ per punctum, cuius abscissa $=x$, ordinata (ab altera lineae abscissarum plaga positive, ab altera negative sumta) $=y$.

Hoc pacto dici potest, quamlibet quantitatem complexam mensurare inaequalitatem inter situm puncti ad quod refertur atque situm puncti initialis, denotante unitate positiva deflexum arbitrarium determinatum versus directionem arbitrariam determinatam; unitate negativa deflexum aeque magnum versus directionem oppositam; denique unitatibus imaginariis deflexus aeque magnos versus duas directiones laterales normales. Hoc modo metaphysica quantitatum, quas imaginarias dicimus, insigniter illustratur.

Si punctum initiale per $(0)$ denotatur, atque duae quantitates complexae $m$, $m'$ ad puncta $M$, $M'$ referuntur, quorum situm relative ad $(0)$ exprimunt, differentia $m-m'$ nihil aliud erit nisi situs puncti $M$ relative ad punctum $M'$: contra, producto $mm'$ repraesentante situm puncti $N$ relative ad $(0)$, facile perspicies, hunc situm perinde determinari per situm puncti $M$ ad $(0)$, ut situs puncti $M'$ determinatur per situm puncti cui respondet unitas positiva, ita ut haud inepte dicas, situs punctorum respondentium quantitatibus complexis $mm'$, $m$, $m'$, $1$ formare proportionem.

Sed uberiorem huius rei tractationem ad aliam occasionem nobis reservamus. Difficultates, quibus theoria quantitatum imaginariarum involuta putatur, ad magnam partem a denominationibus parum idoneis originem traxerunt (quum adeo quidam usi sint nomine absono quantitatum impossibilium). Si, a conceptibus, quos offerunt varietates duarum dimensionum, (quales in maxima puritate conspiciuntur in intuitionibus spatii) profecti, quantitates positivas directas, negativas inversas, imaginarias laterales nuncupavissemus, pro tricis simplicitas, pro caligine claritas successisset.

\art{39}

Quae in art. praec. prolata sunt, ad quantitates complexas continuas referuntur: in arithmetica, quae tantummodo circa numeros integros versatur, schema numerorum complexorum erit systema punctorum aequidistantium et in rectis aequidistantibus ita dispositorum, ut planum infinitum in infinite multa quadrata aequalia dispertiant.

Omnes numeri per numerum complexum datum $a+bi=m$ divisibiles item infinite multa quadrata formabunt, quorum latera $=\sqrt{aa+bb}$ sive areae $=aa+bb$; quadrata posteriora ad priora inclinata erunt, quoties quidem neuter numerorum $a$, $b$ est $=0$.

Cuivis numero per modulum $m$ non divisibili respondebit punctum vel intra tale quadratum situm vel in latere duobus quadratis contiguo; posterior tamen casus locum habere nequit, nisi $a$, $b$ divisorem communem habent: porro patet, numeros secundum modulum $m$ congruos in quadratis suis locos congruentes occupare.

Hinc facile concluditur, si colligantur omnes numeri intra quadratum determinatum siti, nec non omnes qui forte in duobus eius lateribus non oppositis iaceant, denique his adscribatur numerus per $m$ divisibilis, haberi systema completum residuorum incongruorum secundum modulum $m$, i.e. quemvis integrum alicui ex illis et quidem unico tantum congruum esse debere. Nec difficile foret ostendere, horum residuorum multitudinem aequalem esse moduli normae, puta $=aa+bb$. Sed consultum videtur, hoc gravissimum theorema alio modo pure arithmetico demonstrare.

\art{40}

\noindent\textsc{Theorema.} Secundum modulum complexum datum $m=a+bi$, cuius norma $aa+bb=p$, et pro quo $a$, $b$ sunt numeri inter se primi, quilibet integer complexus congruus erit alicui residuo e serie $0$, $1$, $2$, $3\ldots p-1$, et non pluribus.

\medskip
\noindent\textsc{Demonstr.} I. Sint $\alpha$, $\beta$ integri tales qui faciant $\alpha a+\beta b=1$, unde erit
\[
i=\alpha b-\beta a+m(\beta+\alpha i).
\]
Proposito itaque numero integro complexo $A+Bi$, habebimus
\[
A+Bi=A+(\alpha b-\beta a)B+m(\beta B+\alpha B i).
\]
Quare denotando per $h$ residuum minimum positivum numeri $A+(\alpha b-\beta a)B$ secundum modulum $p$, statuendoque
\[
A+(\alpha b-\beta a)B=h+kp=h+m(ak-bki),
\]
erit
\[
A+Bi=h+m(\beta B+ak+(\alpha B-bk)i),
\]
sive $A+Bi\equiv h \pmod{m}$. Q.E.P.

II. Quoties eidem numero complexo duo numeri reales $h$, $h'$ secundum modulum $m$ congrui sunt, etiam inter se congrui erunt. Statuamus itaque $h-h'=m(c+di)$, unde fit
\[
(h-h')(a-bi)=p(c+di),
\]
adeoque
\[
(h-h')a=pc,\quad (h-h')b=-pd
\]
nec non, propter $a\alpha+b\beta=1$,
\[
h-h'=p(c\alpha-d\beta),\quad \text{i.e.}\quad h\equiv h' \pmod{p}.
\]
Quapropter $h$ et $h'$, siquidem sunt inaequales, ambo simul in complexu numerorum $0$, $1$, $2$, $3\ldots p-1$ contenti esse nequeunt. Q.E.S.

\art{41}

\noindent\textsc{Theorema.} Secundum modulum complexum $m=a+bi$, cuius norma $aa+bb=p$, et pro quo $a$, $b$ non sunt inter se primi, sed divisorem communem maximum $\lambda$ habent (quem positive acceptum supponimus), quilibet numerus complexus congruus est residuo $x+yi$ tali, ut $x$ sit aliquis numerorum
\[
0,1,2,3\ldots \frac{p}{\lambda}-1,
\]
atque $y$ aliquis horum $0$, $1$, $2$, $3\ldots \lambda-1$, et quidem unico tantum inter omnia $p$ residua, quae tali forma gaudent.

\medskip
\noindent\textsc{Demonstr.} I. Accipiendo integros $\alpha,\beta$ ita, ut fiat $\alpha a+\beta b=\lambda$, erit
\[
\lambda i=\alpha b-\beta a+m(\beta+\alpha i).
\]
Iam sit $A+Bi$ numerus complexus propositus, $y$ residuum minimum positivum ipsius $B$ secundum modulum $\lambda$, atque $x$ residuum minimum positivum ipsius
\[
A+(\alpha b-\beta a)\cdot\frac{B-y}{\lambda}
\]
secundum modulum $\frac{p}{\lambda}$, statuaturque
\[
A+(\alpha b-\beta a)\cdot\frac{B-y}{\lambda}=x+\frac{p}{\lambda}\cdot k.
\]
Hinc erit
\[
\begin{aligned}
A+Bi-(x+yi)
&=\frac{p}{\lambda}\cdot k+(B-y)i-(\alpha b-\beta a)\frac{B-y}{\lambda}\\
&=\frac{p}{\lambda}\cdot k+\frac{B-y}{\lambda}\cdot m(\beta+\alpha i)\\
&=\left(\frac{a}{\lambda}-\frac{b}{\lambda}\,i\right)km+\frac{B-y}{\lambda}(\beta+\alpha i)m,
\end{aligned}
\]
i.e. per $m$ divisibilis, sive $A+Bi\equiv x+yi \pmod{m}$. Q.E.P.

II. Supponamus, secundum modulum $m$ eidem numero complexo congruos esse duos numeros $x+yi$, $x'+y'i$, qui proin etiam inter se congrui erunt secundum modulum $m$. A potiori itaque secundum modulum $\lambda$ congrui erunt, adeoque $y\equiv y' \pmod{\lambda}$. Quodsi igitur uterque $y$, $y'$ inter numeros $0$, $1$, $2$, $3\ldots \lambda-1$ contentus esse supponitur, necessario debet esse $y=y'$.

Hoc pacto vero etiam fiet $x\equiv x' \pmod{m}$, i.e. $x-x'$ per $m$, adeoque $\frac{x-x'}{\lambda}$ integer per $\frac{a}{\lambda}+\frac{b}{\lambda}\,i$ divisibilis, sive
\[
\frac{x-x'}{\lambda}\equiv 0\pmod{\frac{a}{\lambda}+\frac{b}{\lambda}\,i}.
\]
Hinc autem, quum $\frac{a}{\lambda}$, $\frac{b}{\lambda}$ sint numeri inter se primi, concluditur per partem secundam theorematis art. praec., $\frac{x-x'}{\lambda}$ etiam per normam numeri $\frac{a}{\lambda}+\frac{b}{\lambda}\,i$, i.e. per numerum $\frac{p}{\lambda\lambda}$ divisibilem fore, adeoque $x-x'$ per $\frac{p}{\lambda}$. Quapropter si etiam uterque $x$, $x'$ in complexu numerorum $0$, $1$, $2$, $3\ldots \frac{p}{\lambda}-1$ contentus esse supponitur, necessario erit $x=x'$, sive residua $x+yi$, $x'+y'i$ identica. Q.E.S.

Ceterum sponte patet, huc quoque referendum esse casum, ubi modulus est numerus realis, puta $b=0$, et proin $\lambda=\pm a$, nec non eum, ubi modulus est numerus pure imaginarius, puta $a=0$, et proin $\lambda=\pm b$. In utroque casu habetur $\frac{p}{\lambda}=\lambda$.

\art{42}
Referendo itaque omnes numeros complexos secundum modulum datum inter se congruos ad eandem classem, incongruos ad diversas, omnino aderunt $p$ classes totum numerorum integrorum ambitum exhaurientes, denotante $p$ normam moduli. Complexus totidem numerorum e singulis classibus desumtorum exhibebit systema completum residuorum incongruorum, quale in artt. 40, 41 assignavimus.

Et in hocce quidem systemate electio residuorum classes suas quasi repraesentantium innixa erat principio ei, ut in quavis classe adoptaretur residuum $x+yi$ tale, pro quo $y$ habeat valorem minimum, atque inter omnia, quibus idem valor minimus ipsius $y$ inest, id, pro quo valor ipsius $x$ est minimus, exclusis valoribus negativis tum pro $x$ tum pro $y$.

Sed ad alia proposita aliis principiis uti conveniet, imprimisque notandus est modus is, ubi residua talia adoptantur, quae per modulum divisa offerunt quotientes simplicissimos. Manifesto si
\[
\alpha+\beta i,\quad \alpha'+\beta'i,\quad \alpha''+\beta''i,\ldots
\]
sunt quotientes e divisione numerorum congruorum per modulum oriundi, differentiae tum quantitatum $\alpha,\alpha',\alpha'',\ldots$ inter se erunt numeri integri, tum differentiae inter quantitates $\beta,\beta',\beta'',\ldots$. Patetque semper adesse residuum unum, pro quo $\alpha$ et $\beta$ iaceant inter limites $0$ et $1$, limite priori incluso, posteriori excluso: tale residuum simpliciter vocamus \emph{residuum minimum}.

Si magis placet, loco illorum limitum etiam hi adoptari possunt $-\frac12$ et $+\frac12$ (altero admisso, altero excluso): residuum tali limitationi respondens \emph{absolute minimum} dicemus. Circa haec residua minima offerunt se problemata sequentia.

\art{43}
Residuum minimum numeri complexi dati $A+Bi$ secundum modulum $a+bi$, cuius norma $=p$, invenitur sequenti modo. Si $x+yi$ est residuum minimum quaesitum, erit $(x+yi)(a-bi)$ residuum minimum producti $(A+Bi)(a-bi)$ secundum modulum $(a+bi)(a-bi)$, i.e. secundum modulum $p$.

Statuendo itaque
\[
aA+bB=Fp+f,\qquad aB-bA=Gp+g,
\]
ita ut $f,g$ sint residua minima numerorum $aA+bB$, $aB-bA$ secundum modulum $p$, erit
\[
x+yi=\frac{f+gi}{a-bi},
\]
sive
\[
 x=\frac{af-bg}{p}=A-aF+bG,\qquad
 y=\frac{ag+bf}{p}=B-aG-bF.
\]
Manifesto residua minima $f,g$ vel inter limites $0$ et $p-1$, vel inter hos $-\frac12p$ et $+\frac12p$ accipi debent, prout numeri complexi vel residuum simpliciter minimum vel absolute minimum desideratur.

\art{44}
Constructio systematis completi residuorum minimorum pro modulo dato pluribus modis effici potest. Methodus prima ita procedit, ut primo determinentur limites, intra quos termini reales iacere debent, ac dein pro singulis valoribus intra hos limites sitis assignentur limites partium imaginariarum.

Criterium generale residui minimi $x+yi$ pro modulo $a+bi$ in eo consistit, ut tum $ax+by=\xi$, tum $ay-bx=\eta$ iaceat inter limites $0$ et $aa+bb$, quoties de residuis simpliciter minimis agitur, vel inter limites $-\frac12(aa+bb)$ et $+\frac12(aa+bb)$, quoties residua absolute minima desiderantur, limite altero excluso.

Regulae speciales distinctionem casuum, quos varietas signorum numerorum $a,b$ affert, requirerent, cui tamen evolvendae, quum nulli difficultati obnoxia sit, hic immorari supersedemus: sufficiat, methodi indolem per unicum exemplum exposuisse.

Pro modulo $5+2i$ residua simpliciter minima $x+yi$ ita comparata esse debent, ut tum $5x+2y=\xi$, tum $5y-2x=\eta$ aequetur alicui numerorum $0,1,2,3,\ldots,28$.

Aequatio $29x=5\xi-2\eta$ ostendit, valores positivos ipsius $x$ maiores esse non posse quam $\frac{5\cdot28}{29}$, negativos abstrahendo a signo non maiores quam $\frac{2\cdot28}{29}$. Omnes itaque valores admissibiles ipsius $x$ erunt $-1,0,1,2,3,4$.

Pro $x=-1$ debet esse $2y$ aequalis alicui numerorum $5,6,7,\ldots,33$, atque $5y$ alicui horum $-2,-1,0,1,\ldots,26$; hinc valor minimus ipsius $y$ est $+3$, maximus $+5$.

Tractando perinde valores reliquos ipsius $x$, oritur sequens schema omnium residuorum minimorum:
\[
\begin{array}{r|l}
x&y\\ \hline
-1&3,4,5\\
0&0,1,2,3,4,5\\
+1&1,2,3,4,5,6\\
+2&1,2,3,4,5,6\\
+3&2,3,4,5,6\\
+4&2,3,4
\end{array}
\]

Simili modo pro residuis absolute minimis, $\xi$ et $\eta$ alicui numerorum $-14,-13,-12,\ldots,+14$ aequales esse debent; hinc $29x$ nequit esse extra limites $-7\cdot14$ et $+7\cdot14$, adeoque $x$ alicui numerorum $-3,-2,-1,0,1,2,3$ aequalis esse debet.

Pro $x=-3$ erit $2y=\xi-5x=\xi+15$ alicui numerorum $1,2,3,\ldots,29$ aequalis, $5y=\eta+2x=\eta-6$ autem alicui horum $-20,-19,-18,\ldots,+8$: hinc prodit pro $y$ valor unicus $+1$.

Tractando eodem modo valores reliquos ipsius $x$, habemus schema omnium residuorum absolute minimorum:
\[
\begin{array}{r|l}
x&y\\ \hline
-3&+1\\
-2&-2,-1,0,+1,+2\\
-1&-3,-2,-1,0,+1,+2\\
0&-2,-1,0,+1,+2\\
+1&-2,-1,0,+1,+2,+3\\
+2&-2,-1,0,+1,+2\\
+3&-1
\end{array}
\]

\art{45}
In applicatione methodi secundae duos casus distinguere conveniet.

In casu priori, ubi $a$ et $b$ divisorem communem non habent, fiat
\[
\alpha a+\beta b=1,
\]
sitque $k$ residuum minimum positivum ipsius $\beta a-\alpha b$ secundum modulum $p$. Hinc aequationes identicae
\[
a(\beta a-\alpha b)=\beta p-b(\alpha a+\beta b),\qquad
b(\beta a-\alpha b)=-\alpha p+a(\alpha a+\beta b)
\]
docent, esse
\[
ak\equiv -b,\qquad bk\equiv a\pmod p.
\]
Statuendo itaque ut supra
\[
ax+by=\xi,\qquad ay-bx=\eta,
\]
erit
\[
\eta\equiv k\xi,\qquad \xi\equiv -k\eta\pmod p.
\]

Omnes itaque numeri $\xi+\eta i$, quibus residua simpliciter minima $x+yi$ respondent, habebuntur, dum vel pro $\xi$ deinceps accipiuntur valores $0,1,2,3,\ldots,p-1$, et pro $\eta$ residua minima positiva productorum $k\xi$ secundum modulum $p$, vel ordine alio pro $\eta$ illi valores et pro $\xi$ residua minima productorum $-k\eta$.

E singulis $\xi+\eta i$ dein respondentes $x+yi$ invenientur per formulam
\[
x+yi=\frac{\xi+\eta i}{a-bi}
=\frac{a\xi-b\eta}{p}+\frac{a\eta+b\xi}{p}i.
\]
Ceterum obvium est, $\eta$, dum $\xi$ unitate crescat, vel augmentum $k$ vel decrementum $p-k$ pati, adeoque $x+yi$ vel mutationem
\[
\frac{a-kb}{p}+\frac{ak+b}{p}i
\]
vel hanc
\[
\frac{a-kb}{p}+b+\left(\frac{ak+b}{p}-a\right)i,
\]
quae observatio ad constructionem faciliorem reddendam inservit.

Denique patet, si residua absolute minima $x+yi$ desiderentur, haec praecepta eatenus tantum mutari, quatenus ipsi $\xi$ deinceps tribuendi sint valores inter limites $-\frac12p$ et $+\frac12p$, dum pro $\eta$ accipere oporteat residua absolute minima productorum $k\xi$.

Ecce conspectum residuorum minimorum pro modulo $5+2i$ hoc modo adornatorum.

\begin{center}
\small
\textbf{Residua simpliciter minima.}
\begin{longtable}{L@{\quad}L|L@{\quad}L|L@{\quad}L}
\toprule
\xi+\eta i & x+yi & \xi+\eta i & x+yi & \xi+\eta i & x+yi\\
\midrule
0&0&10+25i&+5i&20+21i&+2+5i\\
1+17i&-1+3i&11+13i&+1+3i&21+9i&+3+3i\\
2+5i&+i&12+i&+2+i&22+26i&+2+6i\\
3+22i&+1+4i&13+18i&+1+4i&23+14i&+3+4i\\
4+10i&+2i&14+6i&+2+2i&24+2i&+4+2i\\
5+27i&-1+5i&15+23i&+1+5i&25+19i&+3+5i\\
6+15i&+3i&16+11i&+2+3i&26+7i&+4+3i\\
7+3i&+1+i&17+28i&+1+6i&27+24i&+3+6i\\
8+20i&+4i&18+16i&+2+4i&28+12i&+4+4i\\
9+8i&+1+2i&19+4i&+3+2i&&\\
\bottomrule
\end{longtable}
\end{center}

\begin{center}
\small
\textbf{Residua absolute minima.}
\begin{longtable}{L@{\quad}L|L@{\quad}L|L@{\quad}L}
\toprule
\xi+\eta i & x+yi & \xi+\eta i & x+yi & \xi+\eta i & x+yi\\
\midrule
-14-6i&-2-2i&-4-10i&-2i&+5-2i&+1\\
-13+11i&-3+i&-3+7i&-1+i&+6-14i&+2-2i\\
-12-i&-2-i&-2-5&-i&+7+3i&+1+i\\
-11-13i&-1-3i&-1+12i&-1+2i&+8-9i&+2-i\\
-10+4i&-2&0&0&+9+8i&+1+2i\\
-9-8i&-1-2i&+1-12i&+1-2i&+10-4i&+2\\
-8+9i&-2+i&+2+5i&+i&+11+13i&+1+3i\\
-7-3i&-1-i&+3-7i&+1-i&+12+i&+2+i\\
-6+14i&-2+2i&+4+10i&+2i&+13-11i&+3-i\\
-5+2i&-1&&&+14+6i&+2+2i\\
\bottomrule
\end{longtable}
\end{center}

Casum secundum, ubi $a,b$ non sunt inter se primi, facile ad casum praecedentem reducere licet. Sit $\lambda$ divisor communis maximus numerorum $a,b$, atque $a=\lambda a'$, $b=\lambda b'$. Denotet $F$ indefinite residuum minimum pro modulo $\lambda$, quatenus tamquam numerus complexus consideratur, i.e. exhibeat indefinite numerum talem $x+yi$ ut $x,y$ sint vel inter limites $0$ et $\lambda$, vel inter hos $-\frac12\lambda$ et $+\frac12\lambda$, prout de residuis vel simpliciter vel absolute minimis agitur. Denotet porro $F'$ indefinite residuum minimum pro modulo $a'+b'i$. Tunc erit
\[
(a'+b'i)F+F'
\]
indefinite residuum minimum pro modulo $a+bi$, prodibitque systema completum horum residuorum, dum omnia $F$ cum omnibus $F'$ combinantur.

\art{46}
Duo numeri complexi inter se primi dicuntur, si praeter unitates alios divisores communes non admittunt: quoties autem tales divisores communes adsunt, ii divisores communes maximi vocantur, quorum norma maxima est.

Si duorum numerorum propositorum resolutio in factores primos praesto est, determinatio divisoris communis maximi prorsus eodem modo perficitur, ut pro numeris realibus (Disquiss. Ar. art. 18). Simul hinc elucet omnes divisores communes duorum numerorum datorum metiri debere eorundem divisorem communem maximum hoc modo inventum.

Quare quum sponte iam pateat, ternos numeros huic socios etiam esse divisores communes, semper quaterni numeri, et non plures, divisores communes maximi appellandi erunt; horumque norma erit multiplum normae cuiusvis alius divisoris communis.

Si resolutio duorum numerorum propositorum in factores simplices non adest, divisor communis maximus adiumento similis algorithmi eruitur, ut pro numeris realibus. Sint $m,m'$ duo numeri propositi, formeturque per divisionem repetitam series $m'',m''',\ldots$, ita ut $m''$ sit residuum absolute minimum ipsius $m$ secundum modulum $m'$, dein $m'''$ residuum absolute minimum ipsius $m'$ secundum modulum $m''$, et sic porro.

Denotando normas numerorum $m,m',m'',m''',\ldots$ respective per $p,p',p'',p''',\ldots$, erit $p''/p'$ norma quotientis $m''/m'$, adeoque per definitionem residui absolute minimi certo non maior quam $\frac12$; idem valet de $p'''/p''$ etc. Quapropter integri reales positivi $p',p'',p''',\ldots$ seriem continuo decrescentem formabunt, unde necessario tandem ad terminum $0$ pervenietur.

Sit hic terminus divisorius $m^{(n+1)}$, statuamusque
\[
\begin{aligned}
 m&=km'+m'',\\
 m'&=k'm''+m''',\\
 m''&=k''m'''+m'''',
\end{aligned}
\]
et sic deinceps, usque ad
\[
m^{(n)}=k^{(n)}m^{(n+1)}.
\]
Percurrendo has aequationes ordine inverso, elucet $m^{(n+1)}$ singulos terminos praecedentes $m^{(n)},\ldots,m'',m',m$ metiri; percurrendo autem easdem ordine directo, manifestum est quemvis divisorem communem numerorum $m,m'$ etiam metiri singulos sequentes. Conclusio prior docet $m^{(n+1)}$ esse divisorem communem numerorum $m,m'$; posterior autem hunc divisorem esse maximum.

Ceterum quoties residuum ultimum $m^{(n+1)}$ alicui quatuor unitatum $1,-1,i,-i$ aequale evadit, hoc indicium erit $m$ et $m'$ inter se primos esse.

\art{47}
Si aequationes articuli praecedentis, omissa ultima, ita combinantur, ut $m'',m''',m'''',\ldots,m^{(n)}$ eliminentur, orietur aequatio talis
\[
m^{(n+1)}=hm+h'm',
\]
ubi $h,h'$ erunt integri, et quidem, si designatione in Disquiss. Ar. art. 27 introducta uti placet,
\[
\begin{aligned}
h&=\pm[k',k'',k''',\ldots,k^{(n-1)}]
=\pm[k^{(n-1)},k^{(n-2)},\ldots,k'',k'],\\
h'&=\mp[k,k',k'',k''',\ldots,k^{(n-1)}]
=\mp[k^{(n-1)},k^{(n-2)},\ldots,k'',k',k],
\end{aligned}
\]
valentibus signis superioribus vel inferioribus, prout $n$ par est vel impar.

Hoc theorema ita enunciamus: divisor communis maximus duorum numerorum complexorum $m,m'$ redigi potest ad formam $hm+h'm'$, ita ut $h,h'$ sint integri.

Manifesto enim hoc non solum de eo divisore communi maximo valet, ad quem algorithmus articuli praecedentis deduxit, sed etiam de tribus illi associatis, pro quibus loco coefficientium $h,h'$ accipere oportebit vel hos $hi,h'i$, vel $-h,-h'$, vel $-hi,-h'i$.

Quoties itaque numeri $m,m'$ inter se primi sunt, satisfieri poterit aequationi
\[
1=hm+h'm'.
\]
Propositi sint e.g. numeri $31+6i=m$, $11-20i=m'$. Hic invenimus
\[
\begin{array}{cl@{\qquad}cl}
k&=i,& m''&=+11-5i,\\
k'&=+1-i,& m'''&=+5-4i,\\
k''&=+2,& m''''&=+1+3i,\\
k'''&=-1-2i,& m'''''&=+i,\\
k''''&=+3-i.&&
\end{array}
\]
atque hinc
\[
[k',k'',k''']=-6-5i,
\qquad
[k,k',k'',k''']=+4-10i,
\]
et proin
\[
m'''''=i=(6+5i)m+(4-10i)m',
\]
nec non
\[
1=(5-6i)m+(-10-4i)m',
\]
quod calculo instituto confirmatur.

\art{48}
Per praecedentia omnia, quae ad theoriam congruentiarum primi gradus in arithmetica numerorum complexorum requiruntur, praeparata sunt: sed quum illa essentialiter non differat ab ea, quae pro arithmetica numerorum realium locum habet, atque in Disquisitionibus Arithmeticis copiose exposita est, praecipua momenta hic adscripsisse sufficiet.

I. Congruentia
\[
mt\equiv 1 \pmod {m'}
\]
aequivalet aequationi indeterminatae
\[
mt+m'u=1.
\]
Si huic satisfit per valores $t=h,u=h'$, illius solutio generaliter exhibetur per
\[
t\equiv h\pmod {m'}.
\]
Conditio autem solubilitatis est, ut modulus $m'$ cum coefficiente $m$ divisorem communem non habeat.

II. Solutio congruentiae
\[
ax+b\equiv c\pmod M
\]
in casu eo, ubi $a,M$ sunt inter se primi, pendet a solutione huius
\[
at\equiv 1\pmod M.
\]
Cui si satisfacit $t=h$, illius solutio generalis continetur in formula
\[
x\equiv (c-b)h\pmod M.
\]

III. Congruentia $ax+b\equiv c\pmod M$ in casu eo, ubi $a,M$ divisorem communem $\lambda$ habent, aequivalet huic
\[
\frac a\lambda x\equiv \frac{c-b}{\lambda}\pmod{M/\lambda}.
\]
Dum itaque pro $\lambda$ adoptatur divisor communis maximus numerorum $a,M$, solutio congruentiae propositae ad casum praecedentem reducitur; patetque ad resolubilitatem requiri et sufficere, ut $\lambda$ etiam differentiam $c-b$ metiatur.

\art{49}
Hactenus elementaria tantum attigimus, quae tamen nexus caussa omittere non licuit. In disquisitionibus altioribus arithmetica numerorum complexorum arithmeticae realium in eo similis est, quod theoremata elegantiora et simpliciora prodeunt, dum tales modulos, qui sunt numeri primi, solos admittimus: revera illorum extensio ad modulos compositos plerumque prolixior quam difficilior est, et laboris potius quam artis. Quapropter in sequentibus imprimis de modulis primis agetur.

\art{50}
Denotante $X$ functionem indeterminatae $x$ talem
\[
Ax^n+Bx^{n-1}+Cx^{n-2}+\text{etc.}+Mx+N,
\]
ubi $n$ est integer realis positivus, $A,B,C,\ldots$ integri reales vel imaginarii, $m$ autem integer complexus, vocabimus hic quoque radicem congruentiae
\[
X\equiv 0\pmod m
\]
quemlibet integrum, qui pro $x$ substitutus ipsi $X$ valorem per modulum $m$ divisibilem conciliat. Solutiones per radices secundum modulum congruas non spectabimus tamquam diversas.

Quoties modulus est numerus primus, talis congruentia ordinis $n$ hic quoque plures quam $n$ solutiones diversas admittere non potest.

Denotante $\alpha$ integrum quemvis determinatum complexum, $X$ adiumento divisionis per $x-\alpha$ indefinite ad formam
\[
X=(x-\alpha)X'+h
\]
reduci potest, ita ut $h$ fiat integer determinatus atque $X'$ functio ordinis $n-1$ cum coefficientibus integris. Iam quoties $\alpha$ est radix congruentiae $X\equiv0\pmod m$, manifesto $h$ divisibilis erit per $m$, sive habebitur indefinite
\[
X\equiv(x-\alpha)X'\pmod m.
\]

Perinde si, denotante $\beta$ integrum determinatum, $X'$ ad formam $(x-\beta)X''+h'$ reducitur, $X''$ erit functio ordinis $n-2$ cum coefficientibus integris. Si vero $\beta$ supponitur esse radix congruentiae $X\equiv0$, etiam satisfacere debet huic $(\beta-\alpha)X'\equiv0$, nec non huic $X'\equiv0$, siquidem radices $\alpha,\beta$ sunt incongruae. Unde colligimus, etiam $h'$ per $m$ divisibilem esse debere, sive indefinite
\[
X\equiv (x-\alpha)(x-\beta)X''\pmod m.
\]

Simili modo accedente radice tertia $\gamma$ prioribus incongrua, habebimus indefinite
\[
X\equiv (x-\alpha)(x-\beta)(x-\gamma)X'''
\]
ita ut $X'''$ sit functio ordinis $n-3$ cum coefficientibus integris. Eodem modo ulterius procedere licet, patetque simul coefficientem termini altissimi in singulis functionibus esse $=A$, quem per $m$ non divisibilem esse supponere licet, alioquin enim congruentia $X\equiv0$ essentialiter ad ordinem inferiorem referenda esset.

Quoties itaque adsunt $n$ radices incongruae, puta $\alpha,\beta,\gamma,\ldots,\nu$, habebimus indefinite
\[
X\equiv A(x-\alpha)(x-\beta)(x-\gamma)\cdots(x-\nu)\pmod m.
\]
Quapropter substitutio novi valoris singulis $\alpha,\beta,\gamma,\ldots,\nu$ incongrui certo ipsi $X$ valorem per $m$ non divisibilem conciliaret, unde theorematis veritas sponte sequitur. Ceterum haec demonstratio essentialiter convenit cum ea, quam in Disq. Ar. art. 43 tradidimus, et cuius singula momenta pro numeris complexis perinde valent ac pro realibus.

\art{51}
Quae in Sectione tertia Disquisitionum Arithmeticarum circa residua potestatum tradita sunt, ad maximam partem, levibus mutationibus adhibitis, etiam in arithmetica numerorum complexorum valent; quin adeo demonstrationes theorematum plerumque retineri possent. Ne tamen quid desit, theoremata principalia demonstrationibus concisis firmata proferemus, ubi semper subintelligendum est modulum esse numerum primum.

\emph{Theorema.} Denotante $k$ integrum per modulum $m$, cuius norma $=p$, non divisibilem, erit
\[
k^{p-1}\equiv1\pmod m.
\]
\emph{Demonstratio.} Constituant $a,b,c,\ldots$ systema completum residuorum incongruorum pro modulo $m$, ita tamen, ut residuum per $m$ divisibile omissum sit, adeoque multitudo illorum numerorum, quorum complexum denotamus per $C$, sit $=p-1$. Sit porro $C'$ complexus productorum $ka,kb,kc,\ldots$.

Ex his productis per hypothesin nullum erit divisibile per $m$, quare singula habebunt residua congrua in complexu $C$, puta fieri poterit
\[
ak\equiv a',\quad bk\equiv b',\quad ck\equiv c',\ldots\pmod m,
\]
ita ut numeri $a',b',c',\ldots$ ipsi in complexu $C$ inveniantur; denotemus complexum numerorum $a',b',c',\ldots$ per $C''$.

Sint $P,P',P''$ producta e singulis numeris complexuum $C,C',C''$ respective, sive
\[
P=abc\cdots,
\qquad
P'=k^{p-1}abc\cdots=k^{p-1}P,
\qquad
P''=a'b'c'\cdots.
\]
Quum numeri complexus $C''$ deinceps congrui sint numeris complexus $C'$, erit $P''\equiv P'$, sive $P''\equiv k^{p-1}P$. At quum facile perspiciatur binos quosvis numeros complexus $C''$ inter se incongruos, adeoque omnes inter se diversos esse, necessario numeri complexus $C''$ cum numeris complexus $C$ prorsus conveniunt, ordine tantummodo mutato; unde fit $P''=P$.

Erit itaque $(k^{p-1}-1)P$ numerus per $m$ divisibilis, unde, quum $m$ sit numerus primus singulos factores ipsius $P$ non metiens, necessario $k^{p-1}-1$ per $m$ divisibilis esse debebit. Q. E. D.

\art{52}
\emph{Theorema.} Denotante $k$, ut in art. praecedenti, integrum per modulum $m$ non divisibilem, atque $t$ exponentem minimum (praeter $0$), pro quo
\[
k^t\equiv1\pmod m,
\]
erit $t$ divisor cuiusvis alius exponentis $u$, pro quo $k^u\equiv1\pmod m$.

\emph{Demonstratio.} Si $t$ non esset divisor ipsius $u$, sit $gt$ multiplum ipsius $t$ proxime maius quam $u$, adeoque $gt-u$ integer positivus minor quam $t$. Ex $k^t\equiv1$, $k^u\equiv1$ sequitur
\[
0\equiv k^{gt}-k^u\equiv k^u(k^{gt-u}-1),
\]
adeoque $k^{gt-u}\equiv1$, i.e. datur potestas ipsius $k$ cum exponente minori quam $t$ unitati congrua, contra hypothesin.

Tamquam corollarium hinc sequitur $t$ certo metiri numerum $p-1$. Numeros tales $k$, pro quibus $t=p-1$, etiam hic radices primitivas pro modulo $m$ vocabimus: quales revera adesse iam ostendemus.

\art{53}
Resolvatur numerus $p-1$ in factores suos primos, ita ut habeatur
\[
p-1=a^\alpha b^\beta c^\gamma\cdots,
\]
designantibus $a,b,c,\ldots$ numeros primos reales positivos inaequales. Sint $A,B,C,\ldots$ integri complexi per $m$ non divisibiles, atque respective congruentiis
\[
x^{(p-1)/a}\equiv1,
\quad
x^{(p-1)/b}\equiv1,
\quad
x^{(p-1)/c}\equiv1,
\quad\text{etc.}
\]
secundum modulum $m$ non satisfacientes, quales dari e theoremate art. 50 manifestum est. Denique sit $h$ congruus secundum modulum $m$ producto
\[
A^{(p-1)/a^\alpha}B^{(p-1)/b^\beta}C^{(p-1)/c^\gamma}\cdots.
\]
Tunc dico $h$ fore radicem primitivam.

\emph{Demonstratio.} Denotando per $t$ exponentem infimae potestatis $h^t$ unitati congruae, erit, si $h$ non esset radix primitiva, $t$ submultiplum ipsius $p-1$, sive $(p-1)/t$ integer unitate maior. Manifesto hic integer factores suos primos reales inter hos $a,b,c,\ldots$ habebit: supponamus itaque, quod licet, $(p-1)/t$ esse divisibilem per $a$, statuamusque $p-1=atu$.

Erit itaque, propter $h^t\equiv1$, etiam $h^{tu}\equiv1$, sive
\[
A^{\frac{p-1}{a^\alpha}\cdot\frac{p-1}{a}}
B^{\frac{p-1}{b^\beta}\cdot\frac{p-1}{a}}
C^{\frac{p-1}{c^\gamma}\cdot\frac{p-1}{a}}
\cdots\equiv1.
\]
At manifesto $(p-1)/(ab^\beta)$ est integer, adeoque
\[
B^{\frac{p-1}{b^\beta}\cdot\frac{p-1}{a}}
=(B^{p-1})^{(p-1)/(ab^\beta)}\equiv1,
\]
perinde etiam factores reliqui, excepta potestate ipsius $A$, unitati congrui erunt. Quapropter esse debet
\[
A^{\frac{p-1}{a^\alpha}\cdot\frac{p-1}{a}}\equiv1.
\]
Iam determinetur integer positivus $\lambda$ talis, ut fiat
\[
\lambda b^\beta c^\gamma\cdots\equiv1\pmod a,
\]
quod fieri poterit, quum numerus primus $a$ ipsum $b^\beta c^\gamma\cdots$ non metiatur; statuaturque
\[
\lambda b^\beta c^\gamma\cdots=1+a\mu.
\]
Manifesto fit
\[
A^{\lambda\cdot\frac{p-1}{a^\alpha}\cdot\frac{p-1}{a}}\equiv1,
\]
sive, quoniam
\[
\lambda\cdot\frac{p-1}{a^\alpha}\cdot\frac{p-1}{a}
=(1+a\mu)\frac{p-1}{a}=(p-1)\mu+\frac{p-1}{a},
\]
habemus
\[
A^{(p-1)\mu}A^{(p-1)/a}\equiv1,
\]
atque hinc, quum sponte sit $A^{(p-1)\mu}\equiv1$, etiam
\[
A^{(p-1)/a}\equiv1,
\]
quod est contra hypothesin. Suppositio itaque $t$ esse submultiplum ipsius $p-1$ consistere nequit, eritque adeo necessario $h$ radix primitiva.


\bigskip
\begin{center}{\large Continuatio: Articuli 54--65}\end{center}

\art{54}
Denotante $h$ radicem primitivam pro modulo $m$, cuius norma $=p$, termini progressionis
\[
1,\ h,\ hh,\ h^3,\ldots,h^{p-2}
\]
inter se incongrui erunt, unde facile colligitur, quemlibet integrum non divisibilem per modulum uni ex istis congruum esse debere, sive illam seriem exhibere systema completum residuorum incongruorum exclusa cifra. Exponens eius potestatis, cui numerus datus congruus est, vocari potest huius \textit{index}, dum $h$ tamquam \textit{basis} consideratur. Ecce quaedam exempla, ubi cuivis indici residuum absolute minimum apposuimus.

\begin{center}\textit{Exemplum primum.}\par
$m=5+4i,\quad p=41,\quad h=1+2i$\end{center}
\begin{center}\small
\begin{longtable}{L@{\quad}L|L@{\quad}L|L@{\quad}L|L@{\quad}L|L@{\quad}L}
\toprule
\text{Ind.} & \text{Residuum} & \text{Ind.} & \text{Residuum} & \text{Ind.} & \text{Residuum} & \text{Ind.} & \text{Residuum} & \text{Ind.} & \text{Residuum}\\
\midrule
0 & +1 & 8 & -4 & 16 & -2+2i & 24 & +2i & 32 & +1+i\\
1 & +1+2i & 9 & -3+i & 17 & -1+2i & 25 & -3i & 33 & +1+3i\\
2 & +1-i & 10 & -i & 18 & +4i & 26 & +2+2i & 34 & +2\\
3 & +3+i & 11 & +2-i & 19 & +1+3i & 27 & +2+i & 35 & -3\\
4 & -2i & 12 & -1-i & 20 & -1 & 28 & +4 & 36 & +2-2i\\
5 & +3i & 13 & +1-3i & 21 & -1-2i & 29 & +3-i & 37 & +1-2i\\
6 & -2-2i & 14 & -2 & 22 & -1+i & 30 & +i & 38 & -4i\\
7 & -2-i & 15 & +3 & 23 & -3-i & 31 & -2+i & 39 & -1-3i\\
\bottomrule
\end{longtable}
\end{center}

\begin{center}\textit{Exemplum secundum.}\par
$m=7,\quad p=49,\quad h=1+2i$\end{center}
\begin{center}\small
\begin{longtable}{L@{\quad}L|L@{\quad}L|L@{\quad}L|L@{\quad}L|L@{\quad}L}
\toprule
\text{Ind.} & \text{Residuum} & \text{Ind.} & \text{Residuum} & \text{Ind.} & \text{Residuum} & \text{Ind.} & \text{Residuum} & \text{Ind.} & \text{Residuum}\\
\midrule
0 & +1 & 10 & -1-i & 20 & +2i & 30 & +2-2i & 40 & +3\\
1 & +1+2i & 11 & +1-3i & 21 & +3+2i & 31 & -1+2i & 41 & +3-i\\
2 & -3-3i & 12 & -i & 22 & -1+i & 32 & +2 & 42 & -2-2i\\
3 & +3-2i & 13 & +2-i & 23 & -3-i & 33 & +2-3i & 43 & +2+i\\
4 & -3i & 14 & -3+3i & 24 & -1 & 34 & +1+i & 44 & -2i\\
5 & -1-3i & 15 & -2-3i & 25 & -1-2i & 35 & -1+3i & 45 & -3-2i\\
6 & -2+2i & 16 & -3 & 26 & +3+3i & 36 & +i & 46 & +1-i\\
7 & +1-2i & 17 & -3+i & 27 & -3+2i & 37 & -2+i & 47 & +3+i\\
8 & -2 & 18 & +2+2i & 28 & +3i & 38 & +3-3i &  & \\
9 & -2+3i & 19 & -2-i & 29 & +1+3i & 39 & +2+3i &  & \\
\bottomrule
\end{longtable}
\end{center}

\art{55}
Adiicimus circa radices primitivas et algorithmum indicum quasdam observationes, demonstrationibus propter facilitatem omissis.

\begin{enumerate}[label=\Roman*.]
\item Indices secundum modulum $p-1$ congrui in systemate dato residuis secundum modulum $m$ congruis respondent et vice versa.
\item Residua, quae respondent indicibus ad $p-1$ primis, etiam sunt radices primitivae et vice versa.
\item Si accepta radice primitiva $h$ pro basi, radicis alius primitivae $h'$ index est $t$, et vice versa $t'$ index ipsius $h$, dum $h'$ pro basi accipitur, erit $tt'\equiv 1\pmod{p-1}$; et si iisdem positis indices cuiusdam alius numeri in his duobus systematibus resp. sunt $u$, $u'$, erit $tu'\equiv u$, $t'u\equiv u'\pmod{p-1}$.
\item Dum numeri $1$, $1+i$ eorumque terni socii (tamquam nimis ieiuni) a modulis nobis considerandis excluduntur, restant numeri primi ii, quos in art.~34 tertio et quarto loco posuimus. Posteriorum normae erunt numeri primi reales formae $4n+1$; priorum normae autem quadrata numerorum primorum realium imparium: in utroque igitur casu $p-1$ per $4$ divisibilis est.
\item Denotando indicem numeri $-1$ per $u$, erit $2u\equiv0\pmod{p-1}$, adeoque vel $u\equiv0$, vel $u\equiv\frac12(p-1)$: at quum index $0$ respondeat residuo $+1$, index numeri $-1$ necessario debet esse $\frac12(p-1)$.
\item Perinde denotando per $u$ indicem numeri $i$, erit $2u\equiv\frac12(p-1)\pmod{p-1}$, adeoque vel $u\equiv\frac14(p-1)$ vel $u\equiv\frac34(p-1)$. Sed hic ambiguitas ab electione radicis primitivae pendet. Scilicet si radice primitiva $h$ pro basi accepta index numeri $i$ est $\frac14(p-1)$, index fiet $\frac34(p-1)$, dum pro basi accipitur $h^\mu$, designante $\mu$ integrum positivum formae $4n+3$ ad $p-1$ primum, e.g. ipsum numerum $p-2$, et vice versa. Quare semissis altera radicum primitivarum conciliat numero $i$ indicem $\frac14(p-1)$, altera indicem $\frac34(p-1)$, manifestoque pro illis basibus $-i$ indicem $\frac34(p-1)$, pro his indicem $\frac14(p-1)$ habebit.
\item Quoties modulus est numerus primus realis positivus formae $4n+3$, puta $=q$, adeoque $p=qq$, indices omnium numerorum realium per $q+1$ divisibiles erunt; denotante enim $t$ indicem numeri realis $k$, erit, propter $k^{q-1}\equiv1\pmod q$, $(q-1)t\equiv0\pmod{qq-1}$, adeoque $\frac{t}{q+1}$ integer. Perinde indices numerorum pure imaginariorum ut $ki$ per $\frac12(q+1)$ divisibiles erunt. Patet itaque, radices primitivas pro talibus modulis inter solos numeros mixtos quaerendas esse.
\item Contra pro modulo $m$, qui est numerus primus complexus mixtus, cuiusque proin norma $p$ est numerus primus realis formae $4n+1$, radices primitivae quaelibet etiam inter numeros reales eligi possunt, inter quos completum adeo systema residuorum incongruorum monstrare licet (art.~40). Manifesto autem quilibet numerus realis, qui est radix primitiva pro modulo complexo $m$, simul erit in arithmetica numerorum realium radix primitiva pro modulo $p$, et vice versa.
\end{enumerate}

\art{56}
Etiamsi theoria residuorum et non-residuorum quadraticorum in arithmetica numerorum complexorum sub ipsa theoria residuorum biquadraticorum contenta sit, tamen antequam ad hanc transeamus, illius theoremata palmaria hic seorsim proferemus: brevitatis vero caussa de solo casu principali, ubi modulus est numerus primus complexus (impar), hic loquemur.

Sit $m$ talis modulus, atque $p$ eius norma. Manifesto quivis integer (per $m$ non divisibilis, quod hic semper subintelligendum) quadrato secundum modulum $m$ congruus fieri vel potest vel non potest, prout illius index, radice aliqua primitiva pro basi accepta, par est vel impar; in casu priori ille integer residuum quadraticum ipsius $m$ dicetur, in posteriori non-residuum. Hinc concluditur, inter $p-1$ numeros qui systema completum residuorum incongruorum (per $m$ non divisibilium) exhibeant, semissem ad residua quadratica, semissem alteram ad non-residua quadratica referri. Cuivis vero alii numero extra illud systema idem character hoc respectu tribuendus est, quo gaudet numerus systematis illi congruus.

Porro ibinde sequitur, productum e duobus residuis quadraticis, nec non productum e duobus non-residuis esse residuum quadraticum; contra productum e residuo quadratico in non-residuum fieri non-residuum; et generaliter productum e quotcunque factoribus esse residuum quadraticum vel non-residuum, prout multitudo non-residuorum inter factores par sit vel impar.

Pro distinguendis residuis quadraticis a non-residuis statim se offert criterium generale sequens:
\[
 k\ \text{est residuum vel non-residuum quadraticum}\quad
 \Longleftrightarrow\quad
 k^{\frac12(p-1)}\equiv 1\ \text{vel}\ -1\pmod m.
\]
Veritas huius theorematis statim inde sequitur, quod, accepta radice primitiva quacunque pro basi, index potestatis $k^{\frac12(p-1)}$ fit vel $\equiv0$ vel $\equiv\frac12(p-1)$, prout index numeri $k$ par est vel impar.

\art{57}
Facile quidem est, pro modulo dato systema residuorum incongruorum completum in duas classes, puta residua et non-residua quadratica distinguere, quo pacto simul omnibus reliquis numeris classes suae sponte assignantur. At longe altioris indaginis est quaestio de criteriis ad distinguendum modulos eos, pro quibus numerus datus est residuum quadraticum, ab iis, pro quibus est non-residuum.

Quod quidem attinet ad unitates reales $+1$ et $-1$, hae in arithmetica numerorum complexorum sunt reapse quadrata, adeoque etiam residua quadratica pro quovis modulo. Aeque facile e criterio art. praec. sequitur, numerum $i$ (et perinde $-i$) esse residuum quadraticum cuiusvis moduli, cuius norma $p$ sit formae $8n+1$, non-residuum vero cuiusvis moduli, cuius norma sit formae $8n+5$.

Quum manifesto nihil intersit, utrum numerus $m$, an aliquis numerorum ipsi associatorum $im$, $-m$, $-im$ pro modulo adoptetur, supponere licebit, modulum esse associatorum primarium (art.~36, II), adeoque statuendo modulum $=a+bi$, esse $a$ imparem, $b$ parem. Quo pacto quum semper sit $aa\equiv1\pmod 8$, $bb$ vero vel $\equiv0$ vel $\equiv4\pmod 8$, prout $b$ sit pariter par vel impariter par, patet numeros $+i$ et $-i$ in casu priori esse residua quadratica moduli, in posteriori non-residua.

\art{58}
Quum diiudicatio characteris numeri compositi, utrum sit residuum quadraticum an non-residuum, pendeat a characteribus factorum, manifesto sufficiet, si evolutionem criteriorum ad distinguendos modulos, pro quibus numerus datus $k$ sit residuum quadraticum, ab iis, pro quibus sit non-residuum, ad tales valores ipsius $k$ limitemus, qui sint numeri primi, insuperque inter associatos primarii.

In qua investigatione inductio protinus theoremata maxime elegantia suppeditat. Incipiamus a numero $1+i$, qui invenitur esse residuum quadraticum modulorum
\[
-1+2i,\\ +3-2i,\\ -5-2i,\\ -1-6i,\\ +5+4i,\\ +5-4i,\\ -7,
+7+2i,
-5+6i,\ldots
\]
non-residuum quadraticum autem sequentium
\[
-1-2i,
-3,
+3+2i,
+1+4i,
+1-4i,
-5+2i,
-1+6i,
+7-2i,
-5-6i,
-3+8i,
-3-8i,
+5+8i,
+5-8i,
+9+4i,
+9-4i,\ldots
\]
Si hunc conspectum, in quo semper e quaternis modulis associatis primarium apposuimus, attente examinamus, facile animadvertimus, modulos $a+bi$ in priori classe omnes esse tales, pro quibus $a+b\equiv+1\pmod 8$ fiat, in posteriori vero tales, pro quibus $a+b\equiv-3\pmod 8$. Manifesto hoc criterium, si loco moduli primarii $m$ adoptamus associatum $-m$, ita immutari debet, ut pro modulis prioris classis sit $a+b\equiv-1$, pro modulis posterioris $\equiv+3\pmod 8$.

Quare, siquidem inductio non fefellerit, generaliter, designante $a+bi$ numerum primum, in quo $a$ impar, $b$ par, $1+i$ fit eius residuum quadraticum vel non-residuum quadraticum, prout $a+b\equiv\pm1$, vel $\equiv\pm3\pmod 8$.

Pro numero $-1-i$ eadem regula valet, quae pro $1+i$. Contra considerando $1-i$ tamquam productum ex $-i$ in $1+i$, manifestum est, numero $1-i$ eundem characterem competere, qui tribuendus sit ipsi $1+i$, quoties $b$ sit pariter par, oppositum autem, quoties $b$ sit impariter par, unde facile colligitur, $1-i$ esse residuum quadraticum numeri primi $a+bi$, quoties sit $a-b\equiv\pm1$, non-residuum autem, quoties habeatur $a-b\equiv\pm3\pmod 8$, semper supponendo, $a$ esse imparem, $b$ parem.

Ceterum haec secunda propositio e priori etiam deduci potest adiumento theorematis generalioris, quod ita enunciamus:

\begin{quote}
In theoria residuorum quadraticorum character numeri $\alpha+\beta i$ respectu moduli $a+bi$ idem est, qui numeri $\alpha-\beta i$ respectu moduli $a-bi$.
\end{quote}
Demonstratio huius theorematis inde petitur, quod uterque modulus eandem normam $p$ habet, atque quoties $(\alpha+\beta i)^{\frac12(p-1)}-1$ per $a+bi$ divisibilis est, etiam $(\alpha-\beta i)^{\frac12(p-1)}-1$ per $a-bi$ divisibilis evadit; quoties autem $(\alpha+\beta i)^{\frac12(p-1)}+1$ per $a+bi$ divisibilis est, etiam $(\alpha-\beta i)^{\frac12(p-1)}+1$ per $a-bi$ divisibilis esse debet.

\art{59}
Progrediamur ad numeros primos impares. Numerum $-1+2i$ invenimus esse residuum quadraticum modulorum
\[
+3+2i,
+1-4i,
-5+2i,
-5-2i,
-1-6i,
+7-2i,
-3+8i,
+5+8i,
+5-8i,
+9+4i,\ldots
\]
non-residuum autem modulorum
\[
-1-2i,
-3,
+3-2i,
+1+4i,
-1+6i,
+5+4i,
+5-4i,
-7,
+7+2i,
-5+6i,
-5-6i,
-3-8i,
+9-4i,\ldots
\]
Reducendo modulos prioris classis ad residua eorum absolute minima secundum modulum $-1+2i$, haec sola invenimus $+1$ et $-1$, puta $+3+2i\equiv-1$, $+1-4i\equiv-1$, $-5+2i\equiv+1$, $-5-2i\equiv-1$, etc. Contra omnes moduli posterioris classis congrui inveniuntur secundum modulum $-1+2i$ vel ipsi $+i$, vel ipsi $-i$.

At numeri $+1$, $-1$ ipsi sunt residua quadratica moduli $-1+2i$, atque $+i$ et $-i$ eiusdem non-residua: quocirca, quatenus inductioni fidem habere licet, prodit theorema: Numerus $-1+2i$ est residuum vel non-residuum quadraticum numeri primi $a+bi$, prout hic est residuum vel non-residuum quadraticum ipsius $-1+2i$, siquidem $a+bi$ est primarius e quaternis associatis, vel potius, si $a$ est impar, $b$ par.

Ceterum ex hoc theoremate sponte sequuntur theoremata analoga circa numeros $+1-2i$, $-1-2i$, $+1+2i$.

\art{60}
Instituendo similem inductionem circa numerum $-3$ vel $+3$, invenimus, utrumque esse residuum quadraticum modulorum
\[
+3+2i,
+3-2i,
-1+6i,
-1-6i,
-7,
-5+6i,
-5-6i,
-3+8i,
-3-8i,
+9+4i,
+9-4i,\ldots
\]
non-residuum vero horum
\[
-1+2i,
-1-2i,
+1+4i,
+1-4i,
-5+2i,
-5-2i,
+5+4i,
+5-4i,
+7+2i,
+7-2i,
+5+8i,
+5-8i,\ldots
\]
Priores secundum modulum $3$ congrui sunt alicui ex his quatuor numeris $+1$, $-1$, $+i$, $-i$; posteriores autem alicui ex his $+1+i$, $+1-i$, $-1+i$, $-1-i$. Illi sunt ipsa residua quadratica numeri $3$, hi non-residua.

Docet itaque haec inductio, numerum primum $a+bi$, supponendo $a$ imparem, $b$ parem, ad numerum $-3$ (nec non ad $+3$) eandem relationem habere, quam hic habet ad illum, quatenus scilicet alter alterius residuum quadraticum sit aut non-residuum.

Extendendo similem inductionem ad alios numeros primos, ubique hanc elegantissimam reciprocitatis legem confirmatam invenimus, deferimurque ad theorema hocce fundamentale circa residua quadratica in arithmetica numerorum complexorum:

\begin{quote}
Denotantibus $a+bi$, $A+Bi$ numeros primos tales, ut $a$, $A$ sint impares, $b$, $B$ pares: erit vel uterque alterius residuum quadraticum, vel uterque alterius non-residuum.
\end{quote}
At non obstante summa theorematis simplicitate, ipsius demonstratio magnis difficultatibus premitur, quibus tamen hic non immoramur, quum theorema ipsum sit tantummodo casus specialis theorematis generalioris, summam theoriae residuorum biquadraticorum quasi exhaurientis. Ad hanc igitur iam transeamus.

\art{61}
Quae in art.~2 prioris commentationis de notione residui et non-residui biquadratici prolata sunt, etiam ad arithmeticam numerorum complexorum extendimus, et perinde ut illic etiam hic disquisitionem ad modulos tales, qui sunt numeri primi, restringimus: simul plerumque tacite subintelligendum erit, modulum ita accipi, ut sit inter associatos primarius, puta $\equiv1$ secundum modulum $2+2i$, nec non numeros, de quorum charactere (quatenus sint residua biquadratica vel non-residua) agitur, per modulum non esse divisibiles.

Pro modulo itaque dato numeri per eum non divisibiles in tres classes dispertiri possent, quarum prima contineret residua biquadratica, secunda non-residua biquadratica ea, quae sunt residua quadratica, tertia non-residua quadratica.

Sed hic quoque praestat, loco tertiae classis binas stabilire, ut omnino habeantur quaternae. Assumta radice quacunque primitiva pro basi, residua biquadratica habebunt indices per $4$ divisibiles sive formae $4n$; non-residua ea, quae sunt residua quadratica, habebunt indices formae $4n+2$; denique non-residuorum quadraticorum indices erunt partim formae $4n+1$, partim formae $4n+3$. Hoc modo classes quaternae quidem oriuntur, at distinctio inter binas posteriores non esset absoluta, sed ab electione radicis primitivae pro basi assumtae dependens; facile enim perspicitur, semissem radicum primitivarum non-residuo quadratico dato conciliare indicem formae $4n+1$, semissem alteram vero indicem formae $4n+3$.

Quam ambiguitatem ut tollamus, supponemus semper talem radicem primitivam adoptari, pro qua index $\frac14(p-1)$ competat numero $+i$ (conf. art.~55, VI). Hoc pacto classificatio oritur, quam concinnius independenter a radicibus primitivis ita enunciare possumus. Classis prima contineat numeros $k$ eos, pro quibus fit $k^{\frac14(p-1)}\equiv1$; hi numeri sunt moduli residua biquadratica. Classis secunda contineat eos, pro quibus $k^{\frac14(p-1)}\equiv i$. Classis tertia eos, pro quibus $k^{\frac14(p-1)}\equiv-1$. Classis quarta denique eos, pro quibus $k^{\frac14(p-1)}\equiv-i$.

Classis tertia comprehendet non-residua biquadratica ea, quae sunt residua quadratica; inter secundam et quartam non-residua quadratica distributa erunt. Numeris harum classium tribuemus resp. \textit{characteres biquadraticos} $0$, $1$, $2$, $3$. Si characterem $\lambda$ numeri $k$ secundum modulum $m$ ita definimus, ut sit exponens eius potestatis ipsius $i$, cui numerus $k^{\frac14(p-1)}$ congruus est, manifesto characteres secundum modulum $4$ congrui pro aequivalentibus habendi sunt. Ceterum haec notio tantisper ad modulos eos limitatur, qui sunt numeri primi: in continuatione harum disquisitionum ostendemus, quomodo etiam modulis compositis adaptari possit.

\art{62}
Quo facilius inductio copiosa circa numerorum characteres adstrui possit, tabulam compendiosam hic adiungimus, cuius auxilio character cuiusvis numeri propositi respectu moduli, cuius norma valorem $157$ non transscendit, levi opera obtinetur, dummodo ad observationes sequentes attendatur.

Quum character numeri compositi aequalis sit (sive secundum modulum $4$ congruus) aggregato characterum singulorum factorum, sufficit, si pro modulo dato characteres numerorum primorum assignare possumus. Porro quum characteres unitatum $-1$, $i$, $-i$ manifesto sint congrui numeris $\frac12(p-1)$, $\frac14(p-1)$, $\frac34(p-1)$ secundum modulum $4$, etiam sufficiet, characteres numerorum inter associatos primariorum exhibuisse. Denique quum moduli secundum modulum $m$ congrui eundem characterem habeant, sufficit, characteres talium numerorum in tabulam recipere, qui continentur in systemate residuorum absolute minimorum.

Praeterea per ratiocinium simile ut in art.~58 demonstratur, si pro modulo $a+bi$ character numeri $A+Bi$ sit $\lambda$, pro modulo $a-bi$ autem $\lambda'$ sit character numeri $A-Bi$, semper esse $\lambda\equiv-\lambda'\pmod 4$, sive $\lambda+\lambda'$ per $4$ divisibilem: quapropter sufficit, in tabulam recipere modulos, in quibus $b$ est vel $0$ vel positivus.

Ita e.g. si quaeritur character numeri $11-6i$ respectu moduli $-5-6i$, substituimus loco horum numerorum hosce $11+6i$, $-5+6i$; dein determinamus (art.~43) residuum absolute minimum numeri $11+6i$ secundum modulum $-5+6i$, quod fit $-1-4i=-1\times(1+4i)$; quare quum pro modulo $-5+6i$ character ipsius $-1$ sit $30$, character numeri $1+4i$ autem, ex tabula, $2$, erit $32$ sive $0$ character numeri $11+6i$ pro modulo $-5+6i$, et proin per observationem ultimam etiam character numeri $11-6i$ pro modulo $-5-6i$.

Perinde si quaeritur character numeri $-5+6i$ respectu moduli $11+6i$, illius residuum absolute minimum $1-5i$ resolvitur in factores $-i$, $1+i$, $3-2i$, quibus respondent characteres $117$, $0$, $1$, unde character quaesitus erit $118$ sive $2$; idem character etiam numero $-5-6i$ respectu moduli $11-6i$ tribuendus est.

\begin{center}\small

\begin{longtable}{L C L}
\toprule
\text{Modulus} & \text{Character} & \text{Numeri}\\
\midrule
-3 & 3 & 1+i\\
+3+2i & 3 & 1+i\\
+1+4i & 1 & -1+2i\\
 & 3 & 1+i\\
-5+2i & 0 & -1-2i\\
 & 1 & 1+i\\
 & 2 & -1+2i\\
-1+6i & 0 & -3\\
 & 1 & 1+i, -1+2i\\
 & 2 & -1-2i\\
+5+4i & 0 & 1+i\\
 & 1 & -3\\
 & 3 & -1+2i, -1-2i\\
-7 & 0 & -3\\
 & 1 & -1+2i, 3-2i\\
 & 2 & 1+i\\
 & 3 & -1-2i, 3+2i\\
+7+2i & 0 & 1+i, 3+2i, 3-2i, 1-4i\\
 & 1 & -3\\
 & 2 & -1-2i, 1+4i\\
 & 3 & -1+2i\\
-5+6i & 0 & 1+i, -3, 3+2i, 3-2i\\
 & 1 & 1-4i\\
 & 2 & 1+4i\\
 & 3 & -1+2i, -1-2i\\
-3+8i & 0 & -1+2i, 3-2i, 1-4i\\
 & 1 & 1+i, 3+2i\\
 & 2 & -3\\
 & 3 & -1-2i, 1+4i, -5+2i\\
+5+8i & 0 & -1-2i\\
 & 1 & -5-2i, -1+6i\\
 & 2 & -1+2i, 3-2i\\
 & 3 & 1+i, -3, 3+2i, 1+4i, 1-4i\\
+9+4i & 0 & -1+2i, 3+2i\\
 & 1 & 1+i, -1-2i, 3-2i\\
 & 2 & -3, 1+4i\\
 & 3 & 1-4i, -5+2i\\
-1+10i & 0 & 1+i, -1+2i, -1-2i, 3+2i\\
 & 1 & -3\\
 & 2 & 3-2i, -5+2i, 5-4i\\
 & 3 & 1+4i, 1-4i\\
+3+10i & 1 & 1+i, -1-2i, 1-4i\\
 & 2 & -3, 3+2i, 1+4i, -5-2i\\
 & 3 & -1+2i, 3-2i\\
-7+8i & 0 & 1+i, -7\\
 & 1 & 3+2i, 3-2i, 1-4i, -5-2i\\
 & 2 & -1-2i, 1+4i, -5+2i, -1-6i\\
 & 3 & -1+2i, -3, -1+6i\\
-11 & 0 & -3\\
 & 1 & 1+i, 3-2i, 1+4i, -5+2i, 5+4i\\
 & 2 & -1+2i, -1-2i\\
 & 3 & 3+2i, 1-4i, -5-2i, 5-4i\\
-11+4i & 0 & 1+i, -1+2i, 3+2i, 5+4i\\
 & 1 & -1-2i, -1+6i\\
 & 2 & -5+2i\\
 & 3 & -3, 3-2i, 1+4i, 1-4i, -5-2i\\
+7+10i & 0 & 1+4i, 1-4i, -1+6i, -1-6i\\
 & 1 & -1+2i, 3+2i, -5+2i\\
 & 2 & 1+i, 3-2i\\
 & 3 & -1-2i, -3, -5-2i\\
+11+6i & 0 & 1+i, -1+2i, -3, 1+4i, 1-4i, -7\\
 & 1 & -1-2i, 3+2i, 3-2i\\
 & 2 & -5-2i, -1+6i, 5-4i\\
 & 3 & -5+2i, 5+4i, 7-2i\\
\bottomrule
\end{longtable}
\end{center}

\art{63}
Operam nunc dabimus, ut criteria communia modulorum, pro quibus numerus primus datus characterem eundem habet, per inductionem detegamus. Modulos semper supponimus primarios inter associatos, puta tales $a+bi$, pro quibus vel $a\equiv1$, $b\equiv0$, vel $a\equiv3$, $b\equiv2\pmod 4$.

Respectu numeri $1+i$, a quo initium facimus, inductionis lex facilius arripitur, si modulos prioris generis (pro quibus $a\equiv1$, $b\equiv0$) a modulis posterioris generis (pro quibus $a\equiv3$, $b\equiv2$) separamus. Adiumento tabulae art. praec. invenimus respondere

\begin{center}\small
\begin{tabular}{C L}
\toprule
\text{characterem} & \text{modulis primi generis}\\
\midrule
0 & 5+4i, -7+8i, -7-8i, -11+4i\\
1 & 1-4i, -3+8i, -3-8i, 9+4i, -11\\
2 & 5-4i, -7, -11-4i\\
3 & -3, 1+4i, 5+8i, 5-8i, 9-4i\\
\bottomrule
\end{tabular}
\end{center}

Si haec septemdecim exempla attente consideramus, in omnibus invenimus characterem
\[
\equiv \frac14(a-b-1) \pmod 4.
\]
Perinde respondet

\begin{center}\small
\begin{tabular}{C L}
\toprule
\text{characterem} & \text{modulis secundi generis}\\
\midrule
0 & 3-2i, -1-6i, 7+2i, -5+6i, -1+10i, 11+6i\\
1 & -5+2i, -1+6i, 7-2i, -1-10i, 3+10i\\
2 & -1+2i, -5-2i, 3-10i, 7+10i\\
3 & -1-2i, 3+2i, -5-6i, 7-10i, 11-6i\\
\bottomrule
\end{tabular}
\end{center}

In omnibus his viginti exemplis, levi attentione adhibita, invenitur character
\[
\equiv \frac14(a-b-5) \pmod 4.
\]
Facile has duas regulas in unam pro utroque modulorum genere valentem contrahere licet, si perpendimus, $\frac14bb$ esse pro modulis prioris generis $\equiv0$, pro modulis posterioris generis $\equiv1\pmod4$. Est itaque character numeri $1+i$ respectu moduli cuiusvis primi inter associatos primarii
\[
\equiv \frac14(a-b-1-bb) \pmod 4.
\]
Obiter hic annotare convenit, quum $(b+1)^2$ semper sit formae $8n+1$, sive $\frac14(2b+bb)$ par, characterem istum semper parem vel imparem fieri, prout $\frac14(a+b-1)$ par sit vel impar, quod quadrat cum regula pro charactere quadratico in art.~58 prolata.

Quum $\frac14(a-b-1)$, $\frac14(a-b+3)$ sint integri, quorum alter par, alter impar, ipsorum productum par erit, sive
\[
\frac18(a-b-1)(a-b+3)\equiv0\pmod4.
\]
Hinc loco expressionis allatae pro charactere biquadratico haec quoque adoptari potest
\[
\frac14(a-b-1-bb)-\frac18(a-b-1)(a-b+3)=\frac18(-aa+2ab-3bb+1),
\]
quae forma eo quoque nomine se commendat, quod non restringitur ad modulos primarios, sed tantummodo supponit $a$ esse imparem, $b$ parem: manifesto enim in hac suppositione vel $a+bi$, vel $-a-bi$ erit numerus inter associatos primarius, valorque istius formulae pro utroque modulo idem.

\art{64}
Proficiscendo a regula ultima in art. praec. eruta invenimus esse
\[
\begin{array}{c|c}
\text{numeri} & \text{characterem }\equiv\\
\hline
-1+i & \frac18(aa+2ab-bb-1)\\
-1-i & \frac18(-aa+2ab+bb+1)\\
+1-i & \frac18(aa+2ab+3bb-1)
\end{array}
\]
Hoc statim inde sequitur, quod character ipsius $i$ est $\frac14(aa+bb-1)$, character ipsius $-1$ autem $\frac12(aa+bb-1)\equiv\frac12bb$, quum $aa-1$ semper sit formae $8n$.

Manifesto hae quatuor regulae, etiamsi hactenus ab inductione mutuatae sint, ita inter se sunt nexae, ut quamprimum unius demonstratio absoluta fuerit, tres reliquae simul sint demonstratae. Vix opus est monere, etiam in his regulis tantummodo supponi $a$ imparem, $b$ parem.

Si formulas ad modulos primarios restrictas adhibere non displicet, hac forma uti possumus. Est
\[
\begin{array}{c|c}
\text{numeri} & \text{characterem }\equiv\\
\hline
-1+i & \frac14(-a-b+1-bb)\\
-1-i & \frac14(a-b-1+bb)\\
+1-i & \frac14(-a-b+1+bb)
\end{array}
\]
Formulae simplicissimae prodeunt, si, ut initio inductionis nostrae feceramus, modulos primi et secundi generis distinguimus. Est scilicet character
\[
\begin{array}{c|c|c}
\text{numeri} & \text{pro modulis primi generis} & \text{pro modulis secundi generis}\\
\hline
-1+i & \frac14(-a-b+1) & \frac14(-a-b-3)\\
-1-i & \frac14(a-b-1) & \frac14(a-b+3)\\
+1-i & \frac14(-a-b+1) & \frac14(-a-b+5)
\end{array}
\]

\art{65}
Pro numero $-1+2i$, ad quem iam progredimur, eandem distinctionem inter modulos $a+bi$ eos, pro quibus $a\equiv1$, $b\equiv0$, atque eos, pro quibus $a\equiv3$, $b\equiv2$, quoque adhibebimus. Tabula art.~62 docet, respectu illius numeri respondere

\begin{center}\small
\begin{tabular}{C L}
\toprule
\text{characterem} & \text{modulis primi generis}\\
\midrule
0 & -3+8i, +5-8i, +9+4i, -11+4i\\
1 & +1+4i, +5-4i, -7, -3-8i\\
2 & +1-4i, +5+8i, -7-8i, -11\\
3 & -3, +5+4i, +9-4i, -7+8i, -11-4i\\
\bottomrule
\end{tabular}
\end{center}

Revocatis singulis his modulis ad residua absolute minima secundum modulum $-1+2i$, animadvertimus, omnes, quibus respondet character $0$, esse $\equiv1$; eos, quibus character $1$ respondet, $\equiv i$; eos, quorum character est $2$, fieri $\equiv-1$; denique omnes, quorum character est $3$, fieri $\equiv-i$. At characteres numerorum $1$, $i$, $-1$, $-i$ pro modulo $-1+2i$ ipsi sunt $0$, $1$, $2$, $3$ resp.; quapropter in omnibus his 17 exemplis character numeri $-1+2i$ respectu moduli prioris generis $a+bi$, cum charactere huius numeri respectu moduli $-1+2i$ identicus est.

Perinde adiumento tabulae invenitur, respondere

\begin{center}\small
\begin{tabular}{C L}
\toprule
\text{characterem} & \text{modulis secundi generis}\\
\midrule
0 & +3+2i, -5-2i, -1+10i, -1-10i, +11+6i\\
1 & +3-2i, -1+6i, -5-6i, +7+10i, +7-10i\\
2 & -5+2i, -1-6i, +7-2i\\
3 & -1-2i, +7+2i, -5+6i, +3+10i, +3-10i, +11-6i\\
\bottomrule
\end{tabular}
\end{center}

Revocatis his modulis ad residua minima secundum modulum $-1+2i$, omnia, quibus resp. characteres $0$, $1$, $2$, $3$ respondent, congrua inveniuntur numeris $-1$, $-i$, $+1$, $+i$; his vero ipsis numeris, si vice versa $-1+2i$ pro modulo adoptatur, competunt characteres $2$, $3$, $0$, $1$ resp. Quapropter in omnibus his 19 exemplis character numeri $-1+2i$ respectu moduli secundi generis duabus unitatibus differt a charactere huius numeri respectu numeri $-1+2i$ pro modulo habiti.

Ceterum nullo negotio perspicitur, prorsus similia respectu numeri $-1-2i$ locum habitura esse.

\art{66}
Pro numero $-3$ distinctionem inter modulos primi generis et secundi omittimus, quum eventus doceat, illam hic superfluam esse. Respondet itaque
\begin{center}\small
\begin{tabular}{c p{5.4in}}
\toprule
character & modulis\\
\midrule
$0$ & $-1+6i,\ -1-6i,\ -7,\ -5+6i,\ -5-6i,\ -11,\ 11+6i,\ 11-6i$\\
$1$ & $-1-2i,\ 1-4i,\ -5+2i,\ 5+4i,\ 7+2i,\ 5-8i,\ -1+10i,\ -7-8i,\ -11-4i,\ 7-10i$\\
$2$ & $3+2i,\ 3-2i,\ -3+8i,\ -3-8i,\ 9+4i,\ 3+10i,\ 3-10i$\\
$3$ & $-1+2i,\ 1+4i,\ -5-2i,\ 5-4i,\ 7-2i,\ 5+8i,\ -1-10i,\ -7+8i,\ -11+4i,\ 7+10i$\\
\bottomrule
\end{tabular}
\end{center}

Revocatis his modulis ad residua minima secundum modulum $3$, videmus, eos, quibus respondet character $0$, esse partim $\equiv1$, partim $\equiv-1$; eos, quorum character est $1$, fieri vel $\equiv1-i$, vel $\equiv-1+i$: eos, quorum character est $2$, fieri vel $\equiv i$, vel $\equiv-i$; denique eos, quibus competit character $3$, esse vel $\equiv1+i$, vel $\equiv-1-i$. Ex hac itaque inductione colligimus, characterem numeri $-3$ pro modulo, qui est numerus primus inter associatos primarius, identicum esse cum charactere huius ipsius numeri, dum $3$, sive, quod eodem redit, $-3$ tamquam modulus consideratur.

\art{67}
Simili inductione circa alios numeros primos instituta, invenimus, numeros $3\pm2i$, $-1\pm6i$, $7\pm2i$, $-5\pm6i$ etc. suppeditare theoremata ei similia, ad quod in art.~65 respectu numeri $-1+2i$ pervenimus; contra numeros $1\pm4i$, $5\pm4i$, $-3\pm8i$, $5\pm8i$, $9\pm4i$ etc. perinde se habere ut numerum $-3$.

Inductio itaque perducit ad elegantissimum theorema, quod ad instar theoriae residuorum quadraticorum in arithmetica numerorum realium \textit{Theorema fundamentale theoriae residuorum biquadraticorum} nuncupare liceat, scilicet:

Denotantibus $a+bi$, $a'+b'i$ numeros primos diversos inter associatos suos primarios, i.e. secundum modulum $2+2i$ unitati congruos, character biquadraticus numeri $a+bi$ respectu moduli $a'+b'i$ identicus erit cum charactere numeri $a'+b'i$ respectu moduli $a+bi$, si vel uterque numerorum $a+bi$, $a'+b'i$, vel alteruter saltem, ad primum genus refertur, i.e. secundum modulum $4$ unitati congruus est: contra characteres illi duabus unitatibus inter se different, si neuter numerorum $a+bi$, $a'+b'i$ ad primum genus refertur, i.e. si uterque secundum modulum $4$ congruus est numero $3+2i$.

At non obstante summa huius theorematis simplicitate, ipsius demonstratio inter mysteria arithmeticae sublimioris maxime recondita referenda est, ita ut, saltem ut nunc res est, per subtilissimas tantummodo investigationes enodari possit, quae limites praesentis commentationis longe transgrederentur. Quamobrem promulgationem huius demonstrationis, nec non evolutionem nexus inter hoc theorema atque ea, quae in initio huius commentationis per inductionem stabilire coeperamus, ad commentationem tertiam nobis reservamus. Coronidis tamen loco iam hic trademus, quae ad demonstrationem theorematum in artt.~63, 64 propositorum requiruntur.

\art{68}
Initium facimus a numeris primis $a+bi$ talibus, pro quibus $b=0$ (tertia specie art.~34), ubi itaque (ut numerus inter associatos primarius sit) $a$ debet esse numerus primus realis negativus formae $-(4n+3)$, pro quo scribemus $-q$, quales sunt $-3$, $-7$, $-11$, $-19$ etc.

Denotando per $\lambda$ characterem numeri $1+i$, illo numero pro modulo accepto, esse debet
\[
i^\lambda \equiv (1+i)^{\frac14(qq-1)}
       \equiv 2^{\frac18(qq-1)}\cdot i^{\frac18(qq-1)}
       \pmod q.
\]
Sed constat, $2$ esse residuum quadraticum, vel non-residuum quadraticum ipsius $q$, prout $q$ sit formae $8n+7$, vel formae $8n+3$, unde colligimus, esse generaliter
\[
2^{\frac12(q-1)}\equiv (-1)^{\frac14(q+1)}
      \equiv i^{\frac12(q+1)}\pmod q,
\]
adeoque evehendo ad potestatem exponentis $\frac14(q+1)$
\[
2^{\frac18(qq-1)}\equiv i^{\frac18(q+1)^2}\pmod q.
\]
Aequatio itaque praecedens hanc formam induit
\[
i^\lambda\equiv i^{\frac18(q+1)^2+\frac18(qq-1)}
      \equiv i^{\frac14(qq+q)}\pmod q,
\]
unde sequitur
\[
\lambda\equiv \frac14(qq+q)
      \equiv \frac14(q+1)^2-\frac14(q+1)\pmod4,
\]
sive quum habeatur $\frac14(q+1)^2\equiv0\pmod4$,
\[
\lambda\equiv-\frac14(q+1)\equiv\frac14(a-1)\pmod4.
\]
Quod est ipsum theorema art.~63 pro casu $b=0$.

\art{69}
Longe vero difficilius absolvuntur moduli $a+bi$ tales, pro quibus non est $b=0$ (numeri quartae speciei art.~34), pluresque disquisitiones erunt praemittendae. Normam $aa+bb$, quae erit numerus primus realis formae $4n+1$, designabimus per $p$.

Denotetur per $S$ complexus omnium residuorum simpliciter minimorum pro modulo $a+bi=m$, exclusa cifra, ita ut multitudo numerorum in $S$ contentorum sit $=p-1$. Designet $x+yi$ indefinite numerum huius systematis, statuaturque
\[
ax+by=\xi,\qquad ay-bx=\eta.
\]
Erunt itaque $\xi$, $\eta$ integri inter limites $0$ et $p$ exclusive contenti: in casu praesente enim, ubi $a$, $b$ inter se primi sunt, formulae art.~45, puta
\[
\eta\equiv k\xi,\qquad \xi\equiv-k\eta\pmod p,
\]
docent, neutrum numerorum $\xi$, $\eta$ esse posse $=0$, nisi alter simul evanescat, adeoque fiat $x=0$, $y=0$, quam combinationem iam eiecimus. Criterium itaque numeri $x+yi$ in $S$ contenti, consistit in eo, ut quatuor numeri $\xi$, $\eta$, $p-\xi$, $p-\eta$ sint positivi.

Praeterea observamus pro nullo tali numero esse posse $\xi=\eta$; hinc enim sequeretur
\[
p(x+y)=a(\xi+\eta)+b(\xi-\eta)=2a\xi,
\]
quod est absurdum, quum nullus factorum $2$, $a$, $\xi$ per $p$ divisibilis sit. Simili ratione aequatio
\[
p(x-y+a+b)=2a\xi+(a+b)(p-\xi-\eta)
\]
docet, esse non posse $\xi+\eta=p$.

Quapropter quum numeri $\xi-\eta$, $p-\xi-\eta$ esse debeant vel positivi vel negativi, hinc petimus subdivisionem systematis $S$ in quatuor complexus $C$, $C'$, $C''$, $C'''$, puta ut coniiciantur
\[
\begin{array}{c|l}
\text{in complexum} & \text{numeri pro quibus}\\
\hline
C & \xi-\eta\ \text{positivus, } p-\xi-\eta\ \text{positivus}\\
C' & \xi-\eta\ \text{positivus, } p-\xi-\eta\ \text{negativus}\\
C'' & \xi-\eta\ \text{negativus, } p-\xi-\eta\ \text{negativus}\\
C''' & \xi-\eta\ \text{negativus, } p-\xi-\eta\ \text{positivus}
\end{array}
\]
Criterium itaque numeri complexus $C$ proprie sextuplex est, puta sex numeri $\xi$, $\eta$, $p-\xi$, $p-\eta$, $\xi-\eta$, $p-\xi-\eta$ positivi esse debent; sed manifesto conditiones 2, 5 et 6 iam sponte implicant reliquas. Similia circa complexus $C'$, $C''$, $C'''$ valent, ita ut criteria completa sint triplicia, puta
\[
\begin{array}{c|l}
\text{pro complexu} & \text{positivi esse debent numeri}\\
\hline
C & \eta,\ \xi-\eta,\ p-\xi-\eta\\
C' & p-\xi,\ \xi-\eta,\ \xi+\eta-p\\
C'' & p-\eta,\ \eta-\xi,\ \xi+\eta-p\\
C''' & \xi,\ \eta-\xi,\ p-\xi-\eta
\end{array}
\]
Ceterum vel nobis non monentibus quisque facile intelliget, in repraesentatione figurata numerorum complexorum (vid. art.~39) numeros systematis $S$ intra quadratum contineri, cuius latera iungant puncta numeros $0$, $a+bi$, $(1+i)(a+bi)$, $i(a+bi)$ repraesentantia, et subdivisionem systematis $S$ respondere partitioni quadrati per rectas diagonales. Sed hocce loco ratiocinationibus pure arithmeticis uti maluimus, illustrationem per intuitionem figuratam lectori perito brevitatis caussa linquentes.

\art{70}
Si quatuor numeri complexi $r=x+yi$, $r'=x'+y'i$, $r''=x''+y''i$, $r'''=x'''+y'''i$ ita inter se nexi sunt, ut habeatur
\[
r'=m+ir,\qquad r''=m+ir'=(1+i)m-r,\qquad
r'''=m+ir''=im-ir,
\]
atque primus $r$ ad complexum $C$ pertinere supponitur, reliqui $r'$, $r''$, $r'''$ resp. ad complexus $C'$, $C''$, $C'''$ pertinebunt.

Statuendo enim $\xi=ax+by$, $\eta=ay-bx$, $\xi'=ax'+by'$, $\eta'=ay'-bx'$, $\xi''=ax''+by''$, $\eta''=ay''-bx''$, $\xi'''=ax'''+by'''$, $\eta'''=ay'''-bx'''$, invenitur
\[
\begin{aligned}
\eta=p-\xi'&=p-\eta''=\xi''',\\
\xi-\eta=\xi'+\eta'-p&=\eta''-\xi''=p-\xi'''-\eta''',\\
p-\xi-\eta=\xi'-\eta'&=\xi''+\eta''-p=\eta'''-\xi''' ,
\end{aligned}
\]
unde adiumento criteriorum theorematis veritas sponte demanat.

Et quum rursus fiat $r=m+ir'''$, facile perspicietur, si $r$ supponatur pertinere ad $C'$, numeros $r'$, $r''$, $r'''$ pertinere resp. ad $C''$, $C'''$, $C$; si ille ad $C''$, hos ad $C'''$, $C$, $C'$; denique si ille ad $C'''$, hos ad $C$, $C'$, $C''$. Simul hinc colligitur, in singulis complexibus $C$, $C'$, $C''$, $C'''$ aeque multos numeros reperiri, puta $\frac14(p-1)$.

\art{71}
\textit{Theorema.} Si denotante $k$ integrum per $m$ non divisibilem singuli numeri complexus $C$ per $k$ multiplicantur, productorumque residuis simpliciter minimis secundum modulum $m$ inter complexus $C$, $C'$, $C''$, $C'''$ distributis, multitudo eorum, quae ad singulos hos complexus pertinent, resp. per $c$, $c'$, $c''$, $c'''$ denotatur: character numeri $k$ respectu moduli $m$ erit
\[
\equiv c'+2c''+3c'''\pmod4.
\]

\textit{Demonstr.} Sint illa $c$ residua minima ad $C$ pertinentia $\alpha$, $\beta$, $\gamma$, $\delta$ etc.; dein $c'$ residua ad $C'$ pertinentia haec $m+i\alpha'$, $m+i\beta'$, $m+i\gamma'$, $m+i\delta'$ etc.; porro $c''$ residua ad $C''$ pertinentia haec $(1+i)m-\alpha''$, $(1+i)m-\beta''$, $(1+i)m-\gamma''$, $(1+i)m-\delta''$ etc.; denique $c'''$ residua ad $C'''$ pertinentia haec $im-i\alpha'''$, $im-i\beta'''$, $im-i\gamma'''$, $im-i\delta'''$ etc.

Iam consideremus quatuor producta, scilicet

\begin{enumerate}[label=\arabic*.]
\item productum ex omnibus $\frac14(p-1)$ numeris complexum $C$ constituentibus;
\item productum productorum, quae e multiplicatione singulorum horum numerorum per $k$ orta erant;
\item productum e residuis minimis horum productorum, puta e numeris $\alpha$, $\beta$, $\gamma$, $\delta$ etc., $m+i\alpha'$, $m+i\beta'$ etc. etc.;
\item productum ex omnibus $c+c'+c''+c'''$ numeris $\alpha$, $\beta$, $\gamma$, $\delta$ etc., $\alpha'$, $\beta'$, $\gamma'$ etc., $\alpha''$, $\beta''$, $\gamma''$ etc., $\alpha'''$, $\beta'''$, $\gamma'''$ etc.
\end{enumerate}

Denotando haec quatuor producta ordine suo per $P$, $P'$, $P''$, $P'''$, manifesto erit
\[
P'=k^{\frac14(p-1)}P,\qquad
P'\equiv P'',\qquad
P''\equiv P'''i^{c'+2c''+3c'''}\pmod m
\]
et proin
\[
Pk^{\frac14(p-1)}\equiv P'''i^{c'+2c''+3c'''}\pmod m.
\]
At facile perspicietur, numeros $\alpha'$, $\beta'$, $\gamma'$ etc., $\alpha''$, $\beta''$, $\gamma''$ etc., $\alpha'''$, $\beta'''$, $\gamma'''$ etc. omnes ad complexum $C$ pertinere, atque tum inter se tum a numeris $\alpha$, $\beta$, $\gamma$, $\delta$ etc. diversos esse, sicuti hi ipsi inter se diversi sint. Omnes itaque hi numeri simul sumti, et abstrahendo ab ordine, prorsus identici esse debent cum omnibus numeris complexum $C$ constituentibus, unde colligimus $P=P'''$, adeoque
\[
Pk^{\frac14(p-1)}\equiv Pi^{c'+2c''+3c'''}\pmod m.
\]
Denique quum singuli factores producti $P$ per $m$ non sint divisibiles, hinc concluditur
\[
k^{\frac14(p-1)}\equiv i^{c'+2c''+3c'''}\pmod m,
\]
unde $c'+2c''+3c'''$ erit character numeri $k$ respectu moduli $m$. Q. E. D.

\art{72}
Quo theorema generale art. praec. ad numerum $1+i$ applicari possit, complexum $C$ denuo in duos complexus minores $G$ et $G'$ subdividere oportet, et quidem referemus in complexum $G$ numeros eos $x+yi$, pro quibus $ax+by=\xi$ minor est quam $\frac12p$, in alterum $G'$ eos, pro quibus $\xi$ est maior quam $\frac12p$; multitudinem numerorum in complexibus $G$, $G'$ contentorum resp. per $g$, $g'$ denotabimus, unde erit $g+g'=\frac14(p-1)$.

Criterium completum numerorum ad $G$ pertinentium itaque erit, ut tres numeri $\eta$, $\xi-\eta$, $p-2\xi$ sint positivi: nam conditio tertia pro complexu $C$, secundum quam $p-\xi-\eta$ positivus esse debet, sub illis implicite iam continetur, quum sit
\[
p-\xi-\eta=(\xi-\eta)+(p-2\xi).
\]
Perinde criterium completum numerorum ad $G'$ pertinentium consistet in valoribus positivis trium numerorum $\eta$, $p-\xi-\eta$, $2\xi-p$.

Hinc facile concluditur, productum cuiusvis numeri complexus $G$ per numerum $1+i$ pertinere ad complexum $C'''$; si enim statuitur
\[
\begin{aligned}
(x+yi)(1+i)&=x'+y'i, \quad
ax'+by'=\xi',\quad ay'-bx'=\eta',\\
\text{invenitur}\quad
\xi'&=\xi-\eta,\qquad
\eta'-\xi'=2\eta,\qquad
p-\xi'-\eta'=p-2\xi,
\end{aligned}
\]
i.e. criterium pro numero $x+yi$ complexui $G$ subdito identicum est cum criterio pro numero $x'+y'i$ ad complexum $C'''$ pertinente. Prorsus simili modo ostenditur, productum cuiusvis numeri complexus $G'$ per $1+i$ pertinere ad complexum $C''$.

Erit itaque, si in art. praec. ipsi $k$ valorem $1+i$ tribuimus, $c=0$, $c'=0$, $c''=g'$, $c'''=g$, et proin character numeri $1+i$ fiet
\[
3g+2g'=\frac12(p-1)+g.
\]
Et quum characteres numerorum $i$, $-1$, sint $\frac14(p-1)$, $\frac12(p-1)$, characteres numerorum $-1+i$, $-1-i$, $1-i$ resp. erunt
\[
\frac34(p-1)+g,\qquad g,\qquad \frac14(p-1)+g.
\]
Totus igitur rei cardo iam in investigatione numeri $g$ vertitur.

\art{73}
Quae in artt.~69--72 exposuimus, proprie independentia sunt a suppositione, $m$ esse numerum primarium: abhinc vero saltem supponemus, $a$ imparem, $b$ parem esse, praetereaque $a$, $b$ et $a-b$ esse numeros positivos.

Ante omnia limites valorum ipsius $x$ in complexu $G$ stabilire oportet. Statuendo
\[
ay-bx=\eta,\qquad (a+b)x-(a-b)y=\zeta,\qquad p-2ax-2by=\theta,
\]
criterium numerorum $x+yi$ ad complexum $G$ pertinentium consistit in tribus conditionibus, ut $\eta$, $\zeta$, $\theta$ sint numeri positivi.

Quum fiat
\[
px=(a-b)\eta+a\zeta,\qquad p(a-2x)=a\theta+2b\eta,
\]
manifestum est, $x$ et $2a-x$ esse debere numeros positivos, sive $x$ alicui numerorum $1$, $2$, $3\ldots\frac12(a-1)$ aequalem.

Porro quum sit
\[
(a-b)\theta=2b\zeta+p(a-b-2x),
\]
patet, quamdiu $x$ minor sit quam $\frac12(a-b)$, conditionem secundam (iuxta quam $\zeta$ positivus esse debet) iam implicare tertiam (quod $\theta$ debet esse positivus); contra quoties $x$ sit maior quam $\frac12(a-b)$, conditionem secundam iam contineri sub tertia. Quamobrem pro valoribus ipsius $x$ his $1$, $2$, $3\ldots\frac12(a-b-1)$ tantummodo prospiciendum est, ut $\eta$ et $\zeta$ positivi evadant, sive ut $y$ maior sit quam $\frac{bx}{a}$ et minor quam $\frac{(a+b)x}{a-b}$: pro valore itaque tali dato ipsius $x$ aderunt numeri $x+yi$ omnino
\[
\gbrack{\frac{(a+b)x}{a-b}}-\gbrack{\frac{bx}{a}},
\]
si uncis in eadem significatione utimur, qua iam alibi passim usi sumus.

Contra pro valoribus ipsius $x$ his $\frac12(a-b+1)$, $\frac12(a-b+3)\ldots\frac12(a-1)$ sufficiet, ut ipsis $\eta$ et $\theta$ valores positivi concilientur, sive ut $y$ maior sit quam $\frac{bx}{a}$ et minor quam
\[
\frac{p-2ax}{2b}=\frac12b+\frac{aa-2ax}{2b}.
\]
Quare pro valore tali dato ipsius $x$ aderunt numeri $x+yi$ omnino
\[
\gbrack{\frac12b+\frac{aa-2ax}{2b}}-\gbrack{\frac{bx}{a}}.
\]
Hinc itaque colligimus, multitudinem numerorum complexus $G$ esse
\[
g=\Sigma\gbrack{\frac{(a+b)x}{a-b}}
 +\Sigma\gbrack{\frac12b+\frac{aa-2ax}{2b}}
 -\Sigma\gbrack{\frac{bx}{a}},
\]
ubi in termino primo summatio extendenda est per omnes valores integros ipsius $x$ ab $1$ usque ad $\frac12(a-b-1)$, in secundo ab $\frac12(a-b+1)$ usque ad $\frac12(a-1)$, in tertio ab $1$ usque ad $\frac12(a-1)$.

Si characteristica $\varphi$ in eadem significatione utimur, ut loco citato, puta ut sit
\[
\varphi(t,u)=\gbrack{\frac{u}{t}}+\gbrack{\frac{2u}{t}}+\gbrack{\frac{3u}{t}}+\cdots+\gbrack{\frac{t'u}{t}},
\]
denotantibus $t$, $u$ numeros positivos quoscunque, atque $t'$ numerum $\gbrack{\frac12t}$, terminus ille primus fit $=\varphi(a-b,a+b)$, tertius $=-\varphi(a,b)$; secundus vero fit
\[
=\frac14bb+\sum\gbrack{\frac{aa-2ax}{2b}}.
\]
Sed fit, scribendo terminos inverso ordine,
\[
\Sigma\gbrack{\frac{aa-2ax}{2b}}
=\gbrack{\frac{a}{2b}}+\gbrack{\frac{3a}{2b}}+\gbrack{\frac{5a}{2b}}+\cdots+\gbrack{\frac{(b-1)a}{2b}}
=\varphi(2b,a)-\varphi(b,a).
\]
Formula itaque nostra sequentem induit formam:
\[
g=\varphi(a-b,a+b)+\varphi(2b,a)-\varphi(a,b)-\varphi(b,a)+\frac14bb.
\]

Consideremus primo terminum $\varphi(a-b,a+b)$, qui protinus transmutatur in
\[
\varphi(a-b,2b)+1+2+3+\cdots+\frac12(a-b-1),
\]
sive in
\[
\varphi(a-b,2b)+\frac18((a-b)^2-1).
\]
Dein quum per theorema generale fiat
\[
\varphi(t,u)+\varphi(u,t)=\gbrack{\frac12t}\cdot\gbrack{\frac12u},
\]
dum $t,u$ sunt integri positivi inter se primi, habemus
\[
\varphi(a-b,2b)=\frac12b(a-b-1)-\varphi(2b,a-b),
\]
adeoque
\[
\varphi(a-b,a+b)=\frac18(aa+2ab-3bb-4b-1)-\varphi(2b,a-b).
\]
Disponamus partes ipsius $\varphi(2b,a-b)$ sequenti modo
\[
\begin{aligned}
&\gbrack{\frac{a-b}{2b}}+\gbrack{\frac{3(a-b)}{2b}}+\gbrack{\frac{5(a-b)}{2b}}+\cdots+\gbrack{\frac{(b-1)(a-b)}{2b}}\\
+{}&\gbrack{\frac{a-b}{b}}+\gbrack{\frac{2(a-b)}{b}}+\gbrack{\frac{3(a-b)}{b}}+\cdots+\gbrack{\frac{\frac12b(a-b)}{b}}.
\end{aligned}
\]
Series secunda manifesto fit
\[
=\varphi(b,a-b)=\varphi(b,a)-1-2-3-\cdots-\frac12b
=\varphi(b,a)-\frac18(bb+2b);
\]
seriem primam ordine terminorum inverso ita exhibemus:
\[
\begin{aligned}
&\gbrack{\frac12(a+1-b)-\frac{a}{2b}}
+\gbrack{\frac12(a+3-b)-\frac{3a}{2b}}\\
&+\gbrack{\frac12(a+5-b)-\frac{5a}{2b}}
+\cdots+\gbrack{\frac12(a-1)-\frac{(b-1)a}{2b}},
\end{aligned}
\]
quae expressio, quum denotante $t$ numerum integrum, $u$ fractum, generaliter sit $[t-u]=t-1-[u]$, mutatur in sequentem
\[
\begin{aligned}
&\frac18b(2a-4-b)-\gbrack{\frac{a}{2b}}-\gbrack{\frac{3a}{2b}}
-\gbrack{\frac{5a}{2b}}-\cdots-\gbrack{\frac{(b-1)a}{2b}}\\
={}&\frac18b(2a-4-b)-\varphi(2b,a)+\varphi(b,a).
\end{aligned}
\]
Hinc fit
\[
\varphi(2b,a-b)=2\varphi(b,a)-\varphi(2b,a)+\frac14b(a-3-b)
\]
et proin
\[
\varphi(a-b,a+b)=\varphi(2b,a)-2\varphi(b,a)+\frac18(aa-bb+2b-1).
\]
Substituendo hunc valorem in formula pro $g$ supra tradita, insuperque
\[
\varphi(a,b)+\varphi(b,a)=\frac14b(a-1),
\]
obtinemus
\[
g=2\varphi(2b,a)-2\varphi(b,a)+\frac18(aa-2ab+bb+4b-1).
\]

\art{74}
Per ratiocinia prorsus similia absolvitur casus is, ubi manentibus $a$, $b$ positivis $a-b$ est negativus sive $b-a$ positivus. Aequationes
\[
p(a-2x)=2b\eta+a\theta,\qquad
p(b-a+2x)=2b\zeta+(b-a)\theta
\]
docent, $\frac12a-x$ atque $x+\frac12(b-a)$ positivos, et proin $x$ alicui numerorum
\[
-\frac12(b-a-1),\ -\frac12(b-a-3),\ -\frac12(b-a-5)\ldots+\frac12(a-1)
\]
aequalem esse debere.

Porro ex aequatione $px+(b-a)\eta=a\zeta$ sequitur, pro valoribus negativis ipsius $x$ conditionem, ex qua $\eta$ debet esse positivus, iam contineri sub conditione, ex qua $\zeta$ debet esse positivus, contrarium vero evenire, quoties ipsi $x$ valor positivus tribuatur. Hinc valores ipsius $y$ pro valore determinato negativo ipsius $x$ inter $\frac{(a+b)x}{a-b}$ et $\frac{p-2ax}{2b}$, contra pro valore positivo ipsius $x$ inter $\frac{bx}{a}$ et $\frac{p-2ax}{2b}$ contenti esse debent: manifesto pro $x=0$ hi limites sunt $0$ et $\frac{p-2ax}{2b}$, valore $y=0$ ipso excluso.

Hinc colligitur
\[
g=-\sum\gbrack{\frac{(a+b)x}{a-b}}
+\sum\gbrack{\frac12b+\frac{aa-2ax}{2b}}
-\sum\gbrack{\frac{bx}{a}},
\]
ubi in termino primo summatio extendenda est per omnes valores negativos ipsius $x$ inde a $-1$ usque ad $-\frac12(b-a-1)$; in secunda per omnes valores ipsius $x$ inde a $-\frac12(b-a-1)$ usque ad $\frac12(a-1)$; in tertia per omnes valores positivos ipsius $x$ inde a $+1$ usque ad $\frac12(a-1)$: hoc pacto e summatione prima prodit $-\varphi(b-a,b+a)$, e secunda perinde ut in art. praec. $\frac14bb+\varphi(2b,a)-\varphi(b,a)$, denique e tertia $-\varphi(a,b)$, sive habetur
\[
g=-\varphi(b-a,b+a)+\varphi(2b,a)-\varphi(b,a)-\varphi(a,b)+\frac14bb.
\]
Iam simili modo ut in art. praec. evolvitur
\[
\begin{aligned}
\varphi(b-a,b+a)
&=\varphi(b-a,2b)-\frac18((b-a)^2-1)\\
&=\frac18(3bb-2ab-aa-4b+1)-\varphi(2b,b-a),
\end{aligned}
\]
nec non
\[
\varphi(2b,b-a)=\varphi(2b,a)-2\varphi(b,a)+\frac14b(b-1-a),
\]
adeoque
\[
\varphi(b-a,b+a)=2\varphi(b,a)-\varphi(2b,a)+\frac18(bb-aa-2b+1),
\]
tandemque
\[
g=2\varphi(2b,a)-2\varphi(b,a)+\frac18(aa-2ab+bb+4b-1).
\]
Evictum est itaque, eandem formulam pro $g$ valere, sive sit $a-b$ positivus sive negativus, dummodo $a$, $b$ sint positivi.

\art{75}
Ut reductionem ulteriorem assequamur, statuemus
\[
\begin{aligned}
L&=\gbrack{\frac{a}{2b}}+\gbrack{\frac{2a}{2b}}+\gbrack{\frac{3a}{2b}}+\cdots+\gbrack{\frac{\frac12ba}{2b}},\\
M&=\gbrack{\frac{(\frac12b+1)a}{2b}}+\gbrack{\frac{(\frac12b+2)a}{2b}}+\gbrack{\frac{(\frac12b+3)a}{2b}}+\cdots+\gbrack{\frac{ba}{2b}},\\
N&=\gbrack{\frac{a+b}{2b}}+\gbrack{\frac{2a+b}{2b}}+\gbrack{\frac{3a+b}{2b}}+\cdots+\gbrack{\frac{\frac12ba+b}{2b}}.
\end{aligned}
\]
Quum facile perspiciatur, haberi generaliter $[u]+\left[u+\frac12\right]=[2u]$, quamcunque quantitatem realem denotet $u$, fit $L+N=\varphi(b,a)$, et quum manifesto sit $L+M=\varphi(2b,a)$, erit
\[
\varphi(2b,a)-\varphi(b,a)=M-N.
\]
Porro autem obvium est, aggregatum termini primi seriei $N$ cum penultimo termino seriei $M$, puta
\[
\gbrack{\frac{a+b}{2b}}+\gbrack{\frac{(b-1)a}{2b}},
\]
fieri $=\frac12(a-1)$, atque eandem summam effici e termino secundo seriei $N$ cum antepenultimo seriei $M$, et sic porro: quare quum etiam terminus ultimus seriei $M$ fiat $=\frac12(a-1)$, ultimus vero terminus seriei $N$ sit
\[
=\gbrack{\frac{a+2}{4}}=\frac14(a\mp1),
\]
valente signo superiori vel inferiori, prout $a$ est formae $4n+1$ vel $4n-1$: erit
\[
M+N=\frac14(a-1)b+\frac14(a\mp1)
\]
et proin
\[
\varphi(2b,a)-\varphi(b,a)=\frac14(a-1)b+\frac14(a\mp1)-2N.
\]
Formula itaque pro $g$ in artt.~73 et 74 inventa, transit in sequentem
\[
g=\frac18((a+b)^2-1)+2n-4N,
\]
statuendo $a\mp1=4n$, ubi $n$ erit integer. Sed quum hinc habeatur $1=16nn-8an+aa$, formula haec etiam sequenti modo exhiberi potest:
\[
g=\frac18(-aa+2ab+bb+1)+4\left(\frac12(a+1)n-nn-N\right).
\]
Quapropter quum $g$ sit character numeri $-1-i$ pro modulo $a+bi$, hic character fit
\[
\equiv \frac18(-aa+2ab+bb+1)\pmod4,
\]
quod est ipsum theorema supra (art.~64) per inductionem erutum, sponteque inde demanant theoremata circa characteres numerorum $1+i$, $1-i$, $-1+i$. Quamobrem haec quatuor theoremata, pro casu eo, ubi $a$ et $b$ sunt positivi, iam rigorose sunt demonstrata.

\art{76}
Si manente $a$ positivo $b$ est negativus, statuatur $b=-b'$, ut fiat $b'$ positivus. Quum iam evictum sit, ita pro modulo $a+b'i$ characterem numeri $-1-i$ esse
\[
\equiv \frac18(-aa+2ab'+b'b'+1)\pmod4,
\]
character numeri $-1+i$ pro modulo $a-b'i$ per theorema in art.~62 prolatum erit
\[
\equiv \frac18(aa-2ab'-b'b'-1),
\]
i.e. character numeri $-1+i$ pro modulo $a+bi$ fit
\[
\equiv \frac18(aa+2ab-bb-1);
\]
hoc vero est ipsum theorema in art.~64 allatum, unde tria reliqua circa characteres numerorum $1+i$, $1-i$, $-1-i$ sponte demanant. Quapropter ista theoremata etiam pro casu, ubi $b$ negativus est, demonstrata sunt, scilicet pro omnibus casibus, ubi $a$ est positivus.

Denique si $a$ est negativus, statuatur $a=-a'$, $b=-b'$. Quum itaque per iam demonstrata character numeri $1+i$ respectu moduli $a'+b'i$ sit
\[
\equiv \frac18(-a'a'+2a'b'-3b'b'+1)\pmod4,
\]
nihilque intersit, utrum numerum $a'+b'i$ an oppositum $-a'-b'i$ moduli loco habeamus; manifesto character numeri $1+i$ respectu moduli $a+bi$ est
\[
\equiv \frac18(-aa+2ab-3bb+1),
\]
et similia valent circa characteres numerorum $1-i$, $-1+i$, $-1-i$. Ex his itaque colligitur, demonstrationem theorematum circa characteres numerorum $1+i$, $1-i$, $-1+i$, $-1-i$ (artt.~63, 64) nulli amplius limitationi obnoxiam esse.

\end{document}
